% !TeX spellcheck = pt_BR
\documentclass[12pt,portuguese]{article}

%----------------------------------------------------------------------
% PACOTES ESSENCIAIS E TRADUÇÃO
%----------------------------------------------------------------------
\usepackage[brazilian]{babel}     % Suporte à língua portuguesa
\usepackage[utf8]{inputenc}       % Codificação de entrada
\usepackage[T1]{fontenc}          % Codificação da fonte
\usepackage{lmodern}              % Fonte escalável

%----------------------------------------------------------------------
% PACOTES MATEMÁTICOS
%----------------------------------------------------------------------
\usepackage{amsmath, amsthm, amssymb, amstext, mathtools}
\usepackage{mathrsfs}             % Para fontes como \mathscr
\usepackage{dsfont}               % Para o indicador de conjunto \mathds{1}
\usepackage[all]{xy}              % Para diagramas comutativos
\usepackage{imakeidx}             % Criação de Lista de Símbolos / Índice

% Desigualdades sempre com símbolos inclinados
\let\leq\leqslant
\let\le\leqslant
\let\ge\geqslant
\let\geq\geqslant

%----------------------------------------------------------------------
% FORMATAÇÃO, LAYOUT E CORES
%----------------------------------------------------------------------
\usepackage{geometry}
\usepackage{varwidth}
\usepackage{graphicx}
\usepackage[pdftex]{xcolor}
\usepackage{titlesec}
\usepackage{appendix}
\usepackage{enumitem}
\usepackage{float}
\usepackage{booktabs}
\usepackage{array}
\usepackage{colortbl}
\usepackage{fontawesome5}
\usepackage{lucide-icons}
\usepackage{caption}
\usepackage{stackengine}
\captionsetup{width=0.8\textwidth}

% Definição de Cores
\definecolor{teal}{RGB}{15,118,110}
\definecolor{text}{RGB}{28,25,23}
\definecolor{muted}{RGB}{87,83,78}
\definecolor{card}{RGB}{255,255,255}
\definecolor{border}{RGB}{231,229,228}

% Definição de margens e espaçamento
\setlength{\topmargin}{-0.7in}
\setlength{\oddsidemargin}{0in}
\setlength{\textheight}{9.55in}
\setlength{\textwidth}{6.5in}
\setlength{\baselineskip}{12mm}

% Formatação dos títulos de seção
\titleformat{\section}{\Large\bfseries}{\thesection.}{0.5em}{}
\titleformat{\subsection}{\large\bfseries}{\thesubsection.}{0.5em}{}
\titleformat{\subsubsection}{\normalfont\itshape}{\thesubsubsection.}{0.6em}{}

%----------------------------------------------------------------------
% TIKZ E GRÁFICOS
%----------------------------------------------------------------------
\usepackage{tikz}
\usepackage{pgfplots}
\pgfplotsset{compat=1.17}
\usetikzlibrary{positioning,calc,fit,nfold,decorations.pathreplacing,calligraphy,tikzmark,patterns,arrows.meta,bending,hobby}

%----------------------------------------------------------------------
% REFERÊNCIAS, LINKS E METADADOS DO PDF
%----------------------------------------------------------------------
\usepackage[pdftex, colorlinks, linkcolor={red!50!black}, 
citecolor={red}, urlcolor={blue!80!black}]{hyperref}

\hypersetup{
pdftitle={Funções Características}, pdfauthor={L. Cioletti},
pdfsubject={Função característica, fórmula de inversão, teorema da
continuidade de Lévy, momentos e derivadas}, pdfkeywords={função característica,
fórmula de inversão, convergência fraca, teorema da continuidade de Lévy,
momentos, MAT-UnB} }

% --- PACOTE E COMANDO PARA A LICENÇA CREATIVE COMMONS ---
\usepackage[type={CC},modifier={by-nc-sa},version={4.0}]{doclicense}
\newcommand\blfootnote[1]{%
\begingroup
\renewcommand\thefootnote{}\footnote{#1}%
\addtocounter{footnote}{-1}%
\endgroup
}

%----------------------------------------------------------------------
% ALIASCNT E AMBIENTES DE TEOREMA (Cleveref)
%----------------------------------------------------------------------
\usepackage{aliascnt}
\theoremstyle{definition}
\newtheorem{theorem}{Teorema}

\newaliascnt{definition}{theorem}
\newtheorem{definition}[definition]{Definição}
\aliascntresetthe{definition}

\newaliascnt{example}{theorem}
\newtheorem{example}[example]{Exemplo}
\aliascntresetthe{example}

\newaliascnt{proposition}{theorem}
\newtheorem{proposition}[proposition]{Proposição}
\aliascntresetthe{proposition}

\newaliascnt{corollary}{theorem}
\newtheorem{corollary}[corollary]{Corolário}
\aliascntresetthe{corollary}

\newaliascnt{lemma}{theorem}
\newtheorem{lemma}[lemma]{Lema}
\aliascntresetthe{lemma}

\newaliascnt{remark}{theorem}
\newtheorem{remark}[remark]{Observação}
\aliascntresetthe{remark}

\newtheoremstyle{plainexercise}
  {3pt}{3pt}{\normalfont}{}{\bfseries}{.}{.5em}
  {\thmnumber{#2}\thmnote{ #3}}
\theoremstyle{plainexercise}
\newtheorem{exercise}{Exercício}

% Versão marcada dos exercícios: a estrela fica pendurada na margem esquerda,
% como nas notas anteriores.
\newtheoremstyle{starredexercise}
  {3pt}{3pt}{\normalfont}{}{\bfseries}{.}{.5em}
  {\llap{\textasteriskcentered\,}\thmnumber{#2}\thmnote{ #3}}
\theoremstyle{starredexercise}
\newtheorem{starredexercise}[exercise]{Exercício}
\theoremstyle{definition}

% Ambientes não numerados usados no texto
\newtheorem*{property}{Propriedade}
\newtheorem*{solution}{Solução}

% Se um enumerate for o primeiro conteúdo de um teorema/exercício (cabeçalho
% ainda pendente, \@inlabel true), o \leavevmode dispara o \everypar do amsthm
% que coloca o cabeçalho na própria linha, antes da lista iniciar.
\makeatletter
\setlist[enumerate]{before={\if@inlabel\leavevmode\fi}}
\makeatother

% Nomes em português para os ambientes
\AtBeginDocument{\renewcommand{\proofname}{\normalfont\textbf{Prova}}}
\renewcommand{\qedsymbol}{$\blacksquare$}

% Cleveref (deve vir depois de hyperref e amsthm)
\usepackage[brazilian,nameinlink]{cleveref}

\crefname{theorem}{teorema}{teoremas}
\Crefname{theorem}{Teorema}{Teoremas}
\crefname{definition}{definição}{definições}
\Crefname{definition}{Definição}{Definições}
\crefname{lemma}{lema}{lemas}
\Crefname{lemma}{Lema}{Lemas}
\crefname{proposition}{proposição}{proposições}
\Crefname{proposition}{Proposição}{Proposições}
\crefname{corollary}{corolário}{corolários}
\Crefname{corollary}{Corolário}{Corolários}
\crefname{example}{exemplo}{exemplos}
\Crefname{example}{Exemplo}{Exemplos}
\crefname{exercise}{exercício}{exercícios}
\Crefname{exercise}{Exercício}{Exercícios}
\crefname{remark}{observação}{observações}
\Crefname{remark}{Observação}{Observações}
\crefname{section}{seção}{seções}
\Crefname{section}{Seção}{Seções}

%----------------------------------------------------------------------
% MACROS ADICIONAIS
%----------------------------------------------------------------------
\let\oldemptyset\emptyset
\let\emptyset\varnothing
\graphicspath{{./}{../arquivos-aux/}}

% Seno na convenção em português
\renewcommand{\sin}{\operatorname{sen}}

%----------------------------------------------------------------------
% DOCUMENTO PRINCIPAL
%----------------------------------------------------------------------
\begin{document}

% --- CABEÇALHO COM O LOGO DA UnB ---
\vspace*{-0.8cm}
\begin{figure}[H]
\includegraphics{logo-mat-unb.pdf}
\end{figure}

% --- TÍTULO, AUTOR E DATA ---
\vspace*{0.3cm}
\begin{center}
{\LARGE\bf Funções Características}
\vspace{1cm}

{\large \bf L. Cioletti}
\vspace{0.3cm}

{\footnotesize Outubro de 2026}
\end{center}

% --- RESUMO (ABSTRACT) ---
\begin{abstract}
Nestas notas, introduzimos a função característica de uma variável aleatória e
mostramos que ela determina completamente a distribuição de $X$. Demonstramos
as propriedades elementares
das funções características, a fórmula de inversão e o teorema da
continuidade de Lévy, que relaciona a convergência fraca de distribuições com a
convergência pontual das funções características correspondentes. Estudamos
também a relação entre o comportamento de $\varphi_{X}$ na origem e os momentos de
$X$. Calculamos explicitamente a função característica das
principais distribuições que aparecerão ao longo do curso. Estas notas incluem também 
uma lista de exercícios. A exposição segue de perto a Seção 3.3 do Capítulo 3 da referência
\cite{MR3930614}.
\end{abstract}

% --- LICENÇA NO RODAPÉ ---
\blfootnote{%
\doclicenseImage[imagewidth=2.6em]\ 
\begingroup
\hypersetup{urlcolor=darkgray}
\href{\doclicenseURL}{%
\doclicenseName - Attribution-NonCommercial-ShareAlike 4.0 International%
}%
\endgroup
}

%----------------------------------------------------------------------
% CONTEÚDO
%----------------------------------------------------------------------
\section{Introdução}\label{sec:introducao}

Nas notas de aula anteriores sobre funções geradoras de momentos, estudamos a
função 
\[ M_X(t)=\mathbb{E}\bigl[e^{tX}\bigr]\] 
e vimos que ela gera os momentos
de $X$ por meio de suas derivadas na origem. Embora a expressão que aparece 
no lado direito da igualdade acima esteja sempre bem-definida, podendo ser $+\infty$,
existem casos em que esta função é finita apenas em $t=0$.
O exemplo mais simples desta observação é dado por uma variável aleatória
com distribuição de Cauchy. Como visto nas notas de aula anterior, neste caso, temos 
$M_X(t)=+\infty$, para todo $t\neq0$. 

A função característica é uma função complexa que retem as informações contidas na 
função geradora de momentos e com a vantagem de não sofre degenerescências como as mencionadas acima.
Esta função é muito parecida com a função geradora de momentos e obtida por uma modificação simples
no expoente. Trocamos o número real $t$ pelo número imaginário puro $it$. Este simples
procedimento permite definir, para qualquer variável aleatória $X$, 
uma função a valores complexos $\varphi_{X}:\mathbb{R}\to\mathbb{C}$ da seguinte maneira
\[
\varphi_{X}(t)
=
\mathbb{E}\bigl[e^{itX}\bigr] 
\equiv 
\mathbb{E}\bigl[\cos(tX)\bigr]+ i\, \mathbb{E}\bigl[\mathrm{sen}\, (tX)\bigr] .
\]

A definição dada acima pode ser pensada também em termos da fórmula de Euler,
$e^{itX}=\cos(tX)+i\sin(tX)$, e da generalização natural dos conceitos de esperança
e variáveis aleatórias para o contexto dos números complexos, discutidos na \Cref{subsec:va-complexas} 
Neste contexto mais geral, segue da identidade 
$|e^{itX}|=\cos^2(tX)+\sin^2(tX)=1$ que a variável aleatória $e^{itX}$ é
limitada e, portanto, integrável. Logo $\varphi_{X}(t)$ está bem definida para toda
variável aleatória $X$ e para todo $t\in\mathbb{R}$, sem nenhuma hipótese sobre
os momentos de $X$. Esse ganho de generalidade é possibilitado pelo 
fato de permitirmos que a função $\varphi_{X}$ assuma valores no plano complexo. 


O objetivo destas notas é desenvolver a teoria básica das funções
características em três frentes. Primeiro, mostramos que $\varphi_{X}$ determina completamente a
distribuição de $X$ por meio da chamada fórmula de inversão.
Em seguida, relacionamos a convergência em distribuição de uma sequência 
de variáveis aleatórias com a convergência pontual da sequência de suas
respectivas funções características, o que fornece uma das ferramentas mais eficientes para provar
teoremas de convergência. Por fim, estudamos a relação entre o comportamento de
$\varphi_{X}$ em uma vizinhança da origem e os momentos de $X$.

Estas notas estão organizadas da seguinte maneira. 
A \Cref{sec:preliminares} reúne os fatos auxiliares
sobre variáveis aleatórias complexas, derivadas dentro da esperança e a
integral de Dirichlet, usados no restante do texto. A
\Cref{sec:definicao} apresenta a definição e as propriedades elementares. A
\Cref{sec:exemplos} calcula a função característica das principais
distribuições do curso. A \Cref{sec:inversao} é dedicada à fórmula de inversão e
a suas consequências. A \Cref{sec:continuidade} contém o teorema da
continuidade de Lévy, e a \Cref{sec:momentos} trata da relação entre derivadas na
origem e momentos. Os exercícios propostos estão reunidos na
\Cref{sec:exercicios}.

Os resultados destas notas pressupõem os conceitos e as técnicas desenvolvidos
nas notas de aula anteriores sobre espaços de probabilidade, variáveis
aleatórias, esperança, funções de distribuição e modos de convergência. Em
particular, usaremos sem demonstrar a convergência dominada, a desigualdade de
Jensen e as caracterizações usuais da convergência fraca ali estabelecidas.

\section{Preliminares}\label{sec:preliminares}

Esta seção reúne três fatos técnicos que serão usados repetidamente. O leitor
pode consultá-la à medida que as referências apareçam no texto.

\subsection{Variáveis aleatórias complexas}\label{subsec:va-complexas}

Como visto na seção anterior a função $\varphi_{X}$ assume valores complexos e 
assim por questão de conveniência e simplicidade
vale a pena definir o que é uma variável aleatória complexa 
e o como fico a extensão do conceito de esperança neste contexto.

\begin{definition}\label{def:va-complexa}
Dizemos que $Z:\Omega\to\mathbb{C}$ é uma \emph{variável aleatória complexa} se
$Z=X+iY$, onde $X$ e $Y$ são variáveis aleatórias reais. Nesse caso, definimos
\[
\mathbb{E}[Z] \equiv  \mathbb{E}[X]+i\,\mathbb{E}[Y],
\]
sempre que $\mathbb{E}\bigl[|X|\bigr]<\infty$ e $\mathbb{E}\bigl[|Y|\bigr]<\infty$.
\end{definition}

Escrevemos também $\operatorname{Re}Z=X$ e $\operatorname{Im}Z=Y$ para a parte
real e a parte imaginária de $Z$, e $\overline{Z}=X-iY$ para o conjugado.
Recorde que $|Z|=\sqrt{X^2+Y^2}$.

Usando a identificação natural do plano complexo com o espaço euclidiano $\mathbb{R}^2$
podemos notar imediatamente que valem as seguintes desigualdades
\[
|X|\leqslant |Z|\leqslant |X|+|Y|,
\qquad
|Y|\leqslant |Z|\leqslant |X|+|Y|.
\]
Observe que segue destas desigualdades que $Z$ é integrável se, e somente se, suas partes real e imaginária o
são. Em particular, $e^{itX}$ é integrável para todo $t$, pois
$|e^{itX}|=1$. Temos ainda as propriedades usuais: a esperança é
linear sobre $\mathbb{C}$, comuta com a conjugação, isto é,
$\overline{\mathbb{E}[Z]}=\mathbb{E}\bigl[\overline{Z}\bigr]$  e vale o seguinte lema.

\begin{lemma}\label{lem:esperanca-modulo}
Se $Z$ é uma variável aleatória complexa integrável, então
\[
\bigl|\mathbb{E}[Z]\bigr|\leqslant \mathbb{E}\bigl[|Z|\bigr].
\]
\end{lemma}

\begin{proof}
Se $\mathbb{E}[Z]=0$, a desigualdade é imediata. Suponha $\mathbb{E}[Z]\neq0$.
Como $\overline{\mathbb{E}[Z]}$ é uma constante complexa, a linearidade da
esperança permite introduzi-la na esperança e concluir que
\[
\bigl|\mathbb{E}[Z]\bigr|^2
=\mathbb{E}[Z]\,\overline{\mathbb{E}[Z]}
=\mathbb{E}\bigl[Z\,\overline{\mathbb{E}[Z]}\bigr].
\]
Da igualdade acima segue que o número $\mathbb{E}\bigl[Z\,\overline{\mathbb{E}[Z]}\bigr]$ é 
um número real, pois é igual a $\bigl|\mathbb{E}[Z]\bigr|^2$. Desta observação e da desigualdade
$\operatorname{Re}z\leqslant|z|$, para todo $z\in\mathbb{C}$, concluímos que
\[
\bigl|\mathbb{E}[Z]\bigr|^2 
=
\operatorname{Re}\bigl(\bigl|\mathbb{E}[Z]\bigr|^2 \bigr)
=
\operatorname{Re}\bigl( \mathbb{E}\bigl[Z\,\overline{\mathbb{E}[Z]}\bigr]\bigr)
=
\mathbb{E}\bigl[\operatorname{Re}\bigl(Z\,\overline{\mathbb{E}[Z]}\bigr)\bigr]
\leqslant\mathbb{E}\bigl[\bigl|Z\,\overline{\mathbb{E}[Z]}\bigr|\bigr]
=\bigl|\mathbb{E}[Z]\bigr|\,\mathbb{E}\bigl[|Z|\bigr].
\]
Dividindo por $\bigl|\mathbb{E}[Z]\bigr|>0$, obtemos a desigualdade desejada.
\end{proof}

\subsection{Derivando sob o Símbolo de Esperança}\label{subsec:derivacao-integral}

Em vários pontos destas notas vamos precisar calcular derivadas de funções que são da forma
$t\mapsto\mathbb{E}\bigl[g(t,X)\bigr]$ e isto, na maioria das vezes, será
feito trocando a ordem da derivada com a esperança. O resultado a seguir
garante que isto pode ser feito sob certas condições. A prova dele é obtida
usando o Teorema do Valor Médio e o Teorema da Convergência Dominada.

\begin{theorem}[Derivação dentro da esperança]\label{thm:derivacao-integral}
Sejam $I\subseteq\mathbb{R}$ um intervalo aberto, $X$ uma variável aleatória
real e $g:I\times\mathbb{R}\to\mathbb{C}$. Suponha que:
\begin{enumerate}
\item[(i)] para cada $t\in I$, a variável aleatória $g(t,X)$ é integrável;
\item[(ii)] para quase todo $\omega\in\Omega$, a função
$t\mapsto g(t,X(\omega))$ é derivável em $I$;
\item[(iii)] existe uma variável aleatória integrável $Y$ tal que, quase
certamente,
\[
\left| \frac{\partial}{\partial t}\, g(t,X)\right|\leqslant Y,
\qquad \forall t\in I.
\]
\end{enumerate}
Então $F(t)=\mathbb{E}\bigl[g(t,X)\bigr]$ é derivável em $I$ e
\[
F'(t)
=
\frac{d}{dt} \mathbb{E}\left[\, g(t,X)\right]
=
\mathbb{E}\left[ \frac{\partial}{\partial t}\, g(t,X)\right]
\qquad \forall t\in I.
\]
\end{theorem}

\begin{proof}
Basta tratar o caso real, aplicando o argumento separadamente às partes real e
imaginária de $g$. Fixe $t\in I$ e tome uma sequência $h_n\to0$ com $h_n\neq0$
e $t+h_n\in I$. Pela linearidade da esperança,
\[
\frac{F(t+h_n)-F(t)}{h_n}
=
\mathbb{E}\left[
\frac{g(t+h_n,X)-g(t,X)}{h_n}
\right].
\]
Queremos passar o limite $n\to\infty$ para dentro da esperança. Pelo Teorema
da Convergência Dominada, basta verificar que o quociente converge quase
certamente e que possui uma majorante integrável que não dependa de $n$.

A convergência segue da hipótese (ii): para quase todo $\omega\in\Omega$, a
função $t\mapsto g(t,X(\omega))$ é derivável em $I$, logo o quociente tende a
$\frac{\partial}{\partial t}\,g(t,X(\omega))$ quando $n\to\infty$.

Para obter a majorante, aplicamos o Teorema do Valor Médio à mesma função,
agora entre $t$ e $t+h_n$. Para quase todo $\omega\in\Omega$, existe
$\xi_n(\omega)$ entre $t$ e $t+h_n$ tal que
\[
\frac{g(t+h_n,X(\omega))-g(t,X(\omega))}{h_n}
=
\frac{\partial}{\partial t}\,g(\xi_n(\omega),X(\omega)).
\]
Como $\xi_n(\omega)\in I$, a hipótese (iii) garante que o módulo do quociente
é majorado por $Y(\omega)$, para quase todo $\omega\in\Omega$ e todo $n$. Em
particular, a derivada $\frac{\partial}{\partial t}\,g(t,X)$ é integrável, pois
é majorada em módulo por $Y$.

Podemos então aplicar o Teorema da Convergência Dominada para passar o limite
para dentro da esperança:
\begin{align*}
\lim_{n\to\infty}\frac{F(t+h_n)-F(t)}{h_n}
&=
\lim_{n\to\infty}\mathbb{E}\left[
\frac{g(t+h_n,X)-g(t,X)}{h_n}
\right]
\\[0.25cm]
&=
\mathbb{E}\left[
\lim_{n\to\infty}\frac{g(t+h_n,X)-g(t,X)}{h_n}
\right]
\\[0.25cm]
&=
\mathbb{E}\left[
\frac{\partial}{\partial t}\,g(t,X)
\right].
\end{align*}
Como a sequência $(h_n)$ é arbitrária, o limite do quociente de diferenças
existe e não depende dela. Logo, $F$ é derivável em $t$ e
\[
F'(t)
=
\mathbb{E}\left[
\frac{\partial}{\partial t}\,g(t,X)
\right],
\]
como queríamos demonstrar.
\end{proof}

O papel da hipótese (iii) é fornecer uma majorante integrável que não dependa de
$t$. Sem ela, a troca entre derivada e esperança pode falhar, mesmo que todas as
esperanças envolvidas existam.

Quando $X$ tem densidade $f$, a esperança se escreve
\[
\mathbb{E}\bigl[g(t,X)\bigr]=\int_{\mathbb{R}}g(t,x)f(x)\,dx,
\]
de modo que o teorema cobre também as expressões com densidade, desde que
$\partial_t g(t,x)$ seja majorada em módulo por uma função $h$ com
$\int_{\mathbb{R}}h(x)f(x)\,dx<\infty$. Para uma
integral $\int_a^b g(t,x)\,dx$ que não provenha de uma densidade de
probabilidade, como a que aparece no \Cref{lem:dirichlet}, a mesma prova se
aplica, com as hipóteses análogas, desde que exista uma função $h$ integrável
em $(a,b)$ com $\bigl|\partial_t g(t,x)\bigr|\leqslant h(x)$ para todo
$t\in I$ e quase todo $x\in(a,b)$.

\subsection{A Integral de Dirichlet}\label{subsec:dirichlet}

A prova da fórmula de inversão, na \Cref{sec:inversao}, depende do valor da
integral imprópria
\[
\int_0^\infty \frac{\sin x}{x}\,dx
\]
conhecida como integral de Dirichlet. Mais ainda, ela depende de argumentos
semelhantes aos empregados no cálculo da integral de Dirichlet. Por isso,
vamos apresentar em detalhes como fazer o cálculo desta integral imprópria. 
A ideia é que este cálculo sirva de modelo mais simples para
as questões de convergência e de limitação da prova da fórmula de inversão.

Vale lembrar que existem outras maneiras mais simples e curtas de calcular a integral de Dirichlet,
como, por exemplo, por meio do Teorema dos Resíduos da Análise Complexa. A
rota que seguimos neste texto foi escolhida por permitir apresentar
em contexto mais elementar vários dos argumentos
empregados na prova da fórmula de inversão que é um resultado muito mais profundo
e sofisticado.

\begin{lemma}[Integral de Dirichlet]\label{lem:dirichlet}
\[
\lim_{T\to\infty}\int_0^T \frac{\sin x}{x}\,dx=\frac{\pi}{2}.
\]
\end{lemma}

\begin{proof}
Primeiro observamos que embora a integral acima seja imprópria na origem, neste
ponto não há muitas dificuldades, já que a função
$x\longmapsto(\sin x)/x$ é contínua em $(0,\infty)$ e, pelo limite fundamental
$\lim_{x\to0}(\sin x)/x=1$, pode ser estendida continuamente a $x=0$ com o
valor $1$. A dificuldade está no comportamento no infinito. O problema aqui é que apesar do
integrando tender a zero, ele decai apenas com taxa $1/x$, e a integral do seu
módulo diverge. De fato, para cada inteiro $k\geqslant0$,
\[
\int_{k\pi}^{(k+1)\pi}\frac{|\sin x|}{x}\,dx
\geqslant
\frac{1}{(k+1)\pi}\int_{k\pi}^{(k+1)\pi}|\sin x|\,dx
=
\frac{2}{(k+1)\pi},
\]
e a série harmônica diverge. Portanto, se a integral imprópria convergir,
isso ocorrerá pelo cancelamento entre os intervalos em que $\sin x$ é
positivo e aqueles em que ele é negativo. Temos, assim, duas tarefas:
mostrar que o limite existe e calcular o seu valor.

A prova trata essas tarefas separadamente. No que vamos chamar de Etapa 1, mostramos que a
integral imprópria converge e denotamos seu valor por $L$. A ferramenta principal nesta etapa é a
integração por partes, que transforma a oscilação do seno em um integrando
com decaimento mais rápido. Nas Etapas 2 e 3, calculamos $L$ com o auxílio de
um parâmetro $\varepsilon>0$. A ideia é introduzir uma versão regularizada da
integral, denotada por $h(\varepsilon)$, cuja derivada é uma integral
elementar, e depois mostrar que $h(\varepsilon)\to L$, quando
$\varepsilon\to0$.

\smallskip
\noindent\textbf{Etapa 1: a integral imprópria converge.}
A integração por partes é útil aqui porque derivar $1/x$ produz $-1/x^2$,
que decai mais rapidamente, enquanto uma primitiva de $\sin x$ permanece
limitada. Na integração por partes, também vamos escolher 
a primitiva $1-\cos x$, e não $-\cos x$, para que ela se
anule em $x=0$. Isso fará com que o termo de fronteira na origem não
contribua.
Mais precisamente, para $0<\eta<T$, a integração por partes com $u=1/x$ e $dv=\sin x\,dx$,
fornece a seguinte igualdade
\[
\int_\eta^T\frac{\sin x}{x}\,dx
=
\left[\frac{1-\cos x}{x}\right]_\eta^T
+
\int_\eta^T\frac{1-\cos x}{x^2}\,dx .
\]
Vamos agora fazer $\eta\to0$. A função $(1-\cos x)/x^2$ tem limite $1/2$ na
origem e, portanto, pode ser estendida continuamente a $x=0$. Além disso,
$(1-\cos\eta)/\eta\to0$. Logo, a igualdade anterior se torna
\[
\int_0^T\frac{\sin x}{x}\,dx
=
\frac{1-\cos T}{T}
+
\int_0^T\frac{1-\cos x}{x^2}\,dx .
\]
Nesta nova expressão, o denominador $x$ foi substituído por $x^2$. Como
$0\leqslant1-\cos x\leqslant2$, o novo integrando é majorado por $2/x^2$ em
$[1,\infty)$, cuja integral converge. Perto da origem, ele é limitado pela
extensão contínua mencionada acima. Assim, a última integral tem limite
quando $T\to\infty$. Como também $(1-\cos T)/T\to0$, concluímos que o limite
\[
L=\lim_{T\to\infty}\int_0^T\frac{\sin x}{x}\,dx
=\int_0^\infty\frac{1-\cos x}{x^2}\,dx
\]
existe. Note que embora segunda igualdade reescreva $L$ como uma integral absolutamente
convergente, esta representação ainda não fornece o seu valor numérico.

\smallskip
\noindent\textbf{Etapa 2: uma versão com decaimento exponencial.}
Para calcular $L$, precisamos de uma ideia adicional, pois não dispomos de
uma primitiva elementar de $(\sin x)/x$. O denominador $x$ é o principal
obstáculo. Ele desapareceria se pudéssemos derivar o integrando em relação a
um parâmetro de modo a produzir um fator $x$. O fator $e^{-\varepsilon x}$
tem exatamente essa propriedade: para $x>0$,
\[
\frac{\partial}{\partial\varepsilon}
\left(e^{-\varepsilon x}\frac{\sin x}{x}\right)
=
-e^{-\varepsilon x}\sin x .
\]
Além disso, quando $\varepsilon=0$, esse fator vale $1$, e recuperamos
formalmente a integral original. Estas observações motivam a considerar a 
seguinte função auxiliar 
\[
h(\varepsilon)\equiv \int_0^\infty e^{-\varepsilon x}\,\frac{\sin x}{x}\,dx .
\]
Primeiro observamos que a integral que define $h(\varepsilon)$ 
converge absolutamente para cada $\varepsilon>0$ fixado. 
De fato, $|\sin x|\leqslant x$ para
$x\geqslant0$, logo $|(\sin x)/x|\leqslant1$, para $x>0$ e desta observação segue que
\[
\left|e^{-\varepsilon x}\frac{\sin x}{x}\right|
\leqslant e^{-\varepsilon x}
\qquad \text{e}\qquad 
\int_0^\infty e^{-\varepsilon x}\,dx=\frac{1}{\varepsilon}<\infty .
\]
O fator exponencial, portanto, torna a integral mais fácil de controlar,
enquanto a derivação em $\varepsilon$ elimina o denominador $x$.

Para justificar a derivação sob o sinal de integral, fixe
$\varepsilon_0>0$. Para $\varepsilon>\varepsilon_0$ e $x>0$, temos
\[
\left|
\frac{\partial}{\partial\varepsilon}
\left(e^{-\varepsilon x}\frac{\sin x}{x}\right)
\right|
=
e^{-\varepsilon x}|\sin x|
\leqslant
e^{-\varepsilon_0 x}.
\]
Note que a função $x\longmapsto e^{-\varepsilon_0 x}$ é integrável em $(0,\infty)$ e não
depende de $\varepsilon$. Pela observação sobre integrais em $dx$ que
encerra a subseção anterior, a hipótese análoga à condição (iii) do
\Cref{thm:derivacao-integral} está satisfeita. Logo,
\[
h'(\varepsilon)
=
-\int_0^\infty e^{-\varepsilon x}\sin x\,dx .
\]
Como $\varepsilon_0>0$ é arbitrário, essa fórmula vale para cada
$\varepsilon>0$.

A integral que resta pode ser calculada por duas integrações por partes.
Para isso, considere
\[
I(\varepsilon)\equiv\int_0^\infty e^{-\varepsilon x}\sin x\,dx .
\]
Como $|\sin x|\leqslant1$ e $|\cos x|\leqslant1$, todas as integrais que
aparecerão abaixo convergem absolutamente, pois
$\int_0^\infty e^{-\varepsilon x}\,dx<\infty$. Assim, podemos fazer as
integrações por partes no intervalo limitado $[0,R]$ e, ao final, tomar
$R\to\infty$.

\smallskip
\noindent\textbf{Primeira integração por partes.}
Escolhemos
\[
u=\sin x,
\qquad
dv=e^{-\varepsilon x}\,dx,
\qquad\text{de modo que}\qquad
du=\cos x\,dx,
\qquad
v=-\frac{e^{-\varepsilon x}}{\varepsilon}.
\]
Pela fórmula de integração por partes obtemos
\begin{align*}
\int_0^R e^{-\varepsilon x}\sin x\,dx
&=
\left.-\frac{e^{-\varepsilon x}\sin x}{\varepsilon}\right|_0^R
+\frac{1}{\varepsilon}\int_0^R e^{-\varepsilon x}\cos x\,dx
\\[0.2cm]
&=
-\frac{e^{-\varepsilon R}\sin R}{\varepsilon}
+\frac{1}{\varepsilon}\int_0^R e^{-\varepsilon x}\cos x\,dx .
\end{align*}
Na última igualdade, usamos $\sin0=0$, de modo que apenas o extremo $x=R$
contribui para o termo de fronteira. Já que
\[
\left|-\frac{e^{-\varepsilon R}\sin R}{\varepsilon}\right|
\leqslant\frac{e^{-\varepsilon R}}{\varepsilon}
\longrightarrow0, 
\qquad\text{quando }R\to\infty,
\]
concluímos que
\begin{equation}\label{eq:dirichlet-partes-1}
I(\varepsilon)
=
\frac{1}{\varepsilon}\int_0^\infty e^{-\varepsilon x}\cos x\,dx .
\end{equation}

\smallskip
\noindent\textbf{Segunda integração por partes.}
Para calcular a integral do cosseno que aparece em
\eqref{eq:dirichlet-partes-1}, escolhemos agora
\[
u=\cos x,
\qquad
dv=e^{-\varepsilon x}\,dx,
\qquad\text{de modo que}\qquad
du=-\sin x\,dx,
\qquad
v=-\frac{e^{-\varepsilon x}}{\varepsilon}.
\]
Novamente pela fórmula de integração por partes, temos
\begin{align*}
\int_0^R e^{-\varepsilon x}\cos x\,dx
&=
\left.-\frac{e^{-\varepsilon x}\cos x}{\varepsilon}\right|_0^R
-\frac{1}{\varepsilon}\int_0^R e^{-\varepsilon x}\sin x\,dx
\\[0.2cm]
&=
\frac{1-e^{-\varepsilon R}\cos R}{\varepsilon}
-\frac{1}{\varepsilon}\int_0^R e^{-\varepsilon x}\sin x\,dx .
\end{align*}

Como
$e^{-\varepsilon R}\cos R\to0$ quando $R\to\infty$, segue que
\begin{equation}\label{eq:dirichlet-partes-2}
\int_0^\infty e^{-\varepsilon x}\cos x\,dx
=
\frac{1}{\varepsilon}-\frac{I(\varepsilon)}{\varepsilon}.
\end{equation}

\smallskip
Observe que esta segunda integração por partes fez reaparecer a integral $I(\varepsilon)$
que queríamos calcular. Isso permite obter uma equação algébrica para esta quantidade e
substituindo \eqref{eq:dirichlet-partes-2} em
\eqref{eq:dirichlet-partes-1} verificar que
\[
I(\varepsilon)
=
\frac{1}{\varepsilon}
\left(\frac{1}{\varepsilon}-\frac{I(\varepsilon)}{\varepsilon}\right)
=
\frac{1}{\varepsilon^2}-\frac{I(\varepsilon)}{\varepsilon^2}.
\]
Passando o termo que contém $I(\varepsilon)$ para o lado esquerdo, segue que
\[
I(\varepsilon)+\frac{I(\varepsilon)}{\varepsilon^2}
=\frac{1}{\varepsilon^2},
\qquad\text{ou seja,}\qquad
\frac{1+\varepsilon^2}{\varepsilon^2}\,I(\varepsilon)
=\frac{1}{\varepsilon^2}.
\]
Como $\varepsilon>0$, podemos multiplicar ambos os lados por $\varepsilon^2$
e dividir por $1+\varepsilon^2$, obtendo
\[
I(\varepsilon)=\frac{1}{1+\varepsilon^2}.
\]

Finalmente, como $h'(\varepsilon)=-I(\varepsilon)$, concluímos que
\[
h'(\varepsilon)=-\frac{1}{1+\varepsilon^2}.
\]

Para recuperar $h$ a partir de $h'$ basta aplicar o Teorema Fundamental do Cálculo
e determinar a constante de integração. Para isso basta analisar o comportamento de $h(\varepsilon)$,
quando $\varepsilon\to\infty$. Usando novamente $|(\sin x)/x|\leqslant1$, temos que
\[
|h(\varepsilon)|
\leqslant
\int_0^\infty e^{-\varepsilon x}\,dx
=\frac{1}{\varepsilon},
\]
de modo que $h(\varepsilon)\to0$, quando $\varepsilon\to\infty$.  
Pelo teorema fundamental do cálculo, para
$R>\varepsilon$,
\[
h(R)-h(\varepsilon)
=
-\int_\varepsilon^R\frac{ds}{1+s^2}
=
-\bigl(\arctan R-\arctan\varepsilon\bigr).
\]
Fazendo $R\to\infty$ e usando $h(R)\to0$ e
$\arctan R\to\pi/2$, obtemos
\[
h(\varepsilon)=\frac{\pi}{2}-\arctan\varepsilon .
\]
Em particular,
\begin{equation}\label{eq:dirichlet-regularizada}
\lim_{\varepsilon\to0}h(\varepsilon)=\frac{\pi}{2}.
\end{equation}

\smallskip
\noindent\textbf{Etapa 3: passagem ao limite $\varepsilon\to0$.}
Resta relacionar esse limite com $L$. Não podemos simplesmente substituir
$\varepsilon=0$ na definição de $h$, pois ainda é preciso justificar a troca
do limite em $\varepsilon$ com a integral imprópria. Para isso, vamos
reescrever $h(\varepsilon)$ em termos da função
\[
F(x)=\int_0^x\frac{\sin u}{u}\,du,\qquad x\geqslant0 .
\]
A escolha de $F$ é natural porque $F(x)$ é exatamente a integral truncada
cujo limite queremos calcular. Pela Etapa 1, $F(x)\to L$ quando
$x\to\infty$. Além disso, $F$ é limitada, contínua em $[0,\infty)$ e satisfaz $F(0)=0$. 
Portanto, podemos afirmar que existe algum $C>0$ tal que $|F(x)|\leqslant C$, para todo
$x\geqslant0$. Neste caso temos necessariamente que $|L|\leqslant C$.

\medskip 

Uma nova integração por partes permitirá substituir o integrando oscilante
pela função convergente $F$. Em $(0,R)$, tomamos $u=e^{-\varepsilon x}$ e
$dv=(\sin x/x)\,dx=dF(x)$. Como $F(0)=0$, segue da fórmula de integração por partes que
\[
\int_0^R e^{-\varepsilon x}\frac{\sin x}{x}\,dx
=
e^{-\varepsilon R}F(R)
+
\varepsilon\int_0^R e^{-\varepsilon x}F(x)\,dx .
\]
Quando $R\to\infty$, o termo de fronteira tende a zero, pois
$|e^{-\varepsilon R}F(R)|\leqslant Ce^{-\varepsilon R}$. Logo,
segue da definição de $h(\varepsilon)$ e da igualdade acima que 
\[
h(\varepsilon)=\varepsilon\int_0^\infty e^{-\varepsilon x}F(x)\,dx .
\]
Desta igualdade podemos interpretar $h(\varepsilon)$ como uma média ponderada dos valores de $F(x)$ pelos pesos
$\varepsilon e^{-\varepsilon x}$, cuja integral total é
\begin{align}\label{eq-pesos1}
\int_0^\infty\varepsilon e^{-\varepsilon x}\,dx=1 .
\end{align}
Para ver que essa média ponderada converge para $L$, note que para cada $M>0$ fixado a
contribuição dos pesos ao longo do intervalo $(0,M)$ é dada por
\[
\int_0^M\varepsilon e^{-\varepsilon x}\,dx=1-e^{-\varepsilon M},
\]
que tende a $0$ quando $\varepsilon\to0^+$. Portanto, à medida
que $\varepsilon$ tende a zero, a contribuição dos valores de $F$ em
intervalos limitados torna-se desprezível. Mas, por outro lado, os valores de $F(x)$
para $x$ grande ficam arbitrariamente próximos de $L$. Esta observação é formalizada abaixo. 

Fixe $\delta>0$. Como $F(x)\to L$, quando $x\to\infty$, então podemos escolher $M>0$ tal que
$|F(x)-L|\leqslant\delta$, para todo $x\geqslant M$. Usando a igualdade \eqref{eq-pesos1} temos
\[
L=\varepsilon\int_0^\infty e^{-\varepsilon x}L\,dx
\]
e, portanto,
\[
h(\varepsilon)-L
=
\varepsilon\int_0^\infty e^{-\varepsilon x}\bigl(F(x)-L\bigr)\,dx .
\]
Separamos agora a integral em duas partes. Em $(0,M)$, usamos a cota uniforme
$|F(x)-L|\leqslant|F(x)|+|L|\leqslant2C$. Em $(M,\infty)$, usamos a
escolha de $M$, que garante $|F(x)-L|\leqslant\delta$. Assim,
\begin{align*}
|h(\varepsilon)-L|
&\leqslant
2C\,\varepsilon\int_0^M e^{-\varepsilon x}\,dx
+
\delta\,\varepsilon\int_M^\infty e^{-\varepsilon x}\,dx
\\[0.2cm]
&=
2C\bigl(1-e^{-\varepsilon M}\bigr)
+\delta e^{-\varepsilon M}
\\[0.2cm]
&\leqslant
2C\varepsilon M+\delta .
\end{align*}
Na última desigualdade, usamos $1-e^{-y}\leqslant y$ que é válida para todo $y\geqslant0$ e também que
$e^{-\varepsilon M}\leqslant1$. Escolhendo $\varepsilon$ suficientemente
pequeno para que $2C\varepsilon M<\delta$, obtemos
$|h(\varepsilon)-L|<2\delta$. Como $\delta>0$ é arbitrário, concluímos que
\[
\lim_{\varepsilon\to0}h(\varepsilon)=L .
\]
Comparando esse limite com \eqref{eq:dirichlet-regularizada}, obtemos
$L=\pi/2$, isto é,
\[
\lim_{T\to\infty}\int_0^T\frac{\sin x}{x}\,dx=\frac{\pi}{2},
\]
como queríamos demonstrar.
\end{proof}

\section{Função Característica e Propriedades Elementares}\label{sec:definicao}

Nesta seção, introduzimos o objeto central destas notas e estabelecemos algumas de suas
propriedades elementares. 

\subsection{A definição}\label{subsec:fc-definicao}

Na \Cref{def:va-complexa} fixamos o significado de $\mathbb{E}[Z]$ para uma
variável aleatória complexa $Z$, e no \Cref{lem:esperanca-modulo} vimos que o
módulo da esperança não excede a esperança do módulo. Com essas duas
ferramentas podemos definir o objeto central destas notas.

\begin{definition}[Função característica]\label{def:func-caracteristica}
Seja $X$ uma variável aleatória real. A \emph{função característica} de $X$ é a função
$\varphi_{X}\colon\mathbb{R}\to\mathbb{C}$ dada por
\[
  \varphi_{X}(t) \;=\; \mathbb{E}\bigl[e^{itX}\bigr],
  \qquad \forall t\in\mathbb{R}.
\]
\end{definition}

Observe, antes de tudo, que a definição faz sentido para \emph{toda} variável
aleatória, sem qualquer hipótese sobre momentos. De fato, para cada $t$
fixado, a variável aleatória $e^{itX} = \cos tX + i\sin tX$ tem módulo
unitário:
\[
  \bigl|e^{itX}\bigr| \;=\; \sqrt{\cos^2 tX + \sin^2 tX} \;=\; 1 .
\]
Logo $e^{itX}$ é limitada e, portanto, integrável. Além disso, a \Cref{def:va-complexa}
nos diz que basta calcular as esperanças das partes real e imaginária,
que são as variáveis aleatórias limitadas $\cos tX$ e $\sin tX$.

A fórmula de Euler permite reescrever a função característica em termos da esperança de
duas variáveis aleatórias reais.
Como $e^{itX} = \cos tX + i\sin tX$, segue da definição de esperança de uma variável
aleatória complexa que
\begin{equation}\label{eq:fc-partes}
  \varphi_{X}(t) \;=\; \mathbb{E}\bigl[\cos tX\bigr] + i\,\mathbb{E}\bigl[\sin tX\bigr],
  \qquad t\in\mathbb{R},
\end{equation}
ou equivalentemente,
\[
  \operatorname{Re}\varphi_{X}(t) \;=\; \mathbb{E}\bigl[\cos tX\bigr]
  \qquad\text{e}\qquad
  \operatorname{Im}\varphi_{X}(t) \;=\; \mathbb{E}\bigl[\sin tX\bigr].
\]
Assim, as funções $t\longmapsto \operatorname{Re}\varphi_{X}(t)$ e 
$t\longmapsto \operatorname{Im}\varphi_{X}(t)$ são funções
reais de uma variável real, e consequentemente conhecer $\varphi_{X}$ equivale a conhecer esse par de
funções reais. Em particular, $\varphi_{X}(0) = \mathbb{E}[1] = 1$. Aplicando
ainda o \Cref{lem:esperanca-modulo} a $e^{itX}$ e usando que o seu módulo é
sempre igual a $1$, segue que
\[
  \bigl|\varphi_{X}(t)\bigr| \;\leqslant\; \mathbb{E}\bigl[\bigl|e^{itX}\bigr|\bigr] \;=\; 1
  \qquad\text{para todo } t\in\mathbb{R}.
\]

\medskip

Vale a pena registrar desde já o que se ganha e o que se perde com essa
definição em comparação com a função geradora de momentos. 
O que se ganha é a universalidade: $\varphi_{X}$ existe para toda
variável aleatória $X$, sem que se exija nada sobre os momentos da variável aleatória $X$ 
e além do mais, o seu domínio é sempre toda a reta real. Recorde que a função geradora de
momentos exige $\mathbb{E}[e^{sX}]<\infty$ para algum $s\neq 0$, condição que
pode falhar até para distribuições com todos os momentos finitos. O que se
perde é a simplicidade dos valores: $\varphi_{X}(t)$ é em geral um número
complexo, e não um número real não negativo, de modo que não podemos falar em
monotonicidade nem interpretar $\varphi_{X}(t)$ como uma massa ou um peso. A
informação sobre a distribuição de $X$ fica codificada em $\varphi_{X}$ de uma
forma menos direta, mas, como veremos, sem perda alguma.

\subsection{Algumas Propriedades Elementares}\label{subsec:fc-propriedades}

As cinco propriedades da função característica reunidas abaixo
decorrem diretamente da definição e serão usadas em praticamente todos os
argumentos das próximas seções. 

\begin{theorem}\label{thm:propriedades-fc}
Seja $X$ uma variável aleatória real, com função característica $\varphi_X$. Para
todos $t,h,a,b\in\mathbb{R}$ valem:
\begin{enumerate}[label=(\alph*)]
\item $\varphi_X(0)=1$;
\item $\varphi_X(-t)=\overline{\varphi_X(t)}$;
\item $|\varphi_X(t)|=\bigl|\mathbb{E}[e^{itX}]\bigr|\leqslant\mathbb{E}\bigl[|e^{itX}|\bigr]=1$;
\item $|\varphi_X(t+h)-\varphi_X(t)|\leqslant\mathbb{E}\bigl[|e^{ihX}-1|\bigr]$.
Em particular, a função $\varphi_X$ é uniformemente contínua em $\mathbb{R}$;
\item $\mathbb{E}\bigl[e^{it(aX+b)}\bigr]=e^{itb}\varphi_X(at)$.
\end{enumerate}
\end{theorem}

\begin{proof}

\begin{enumerate}[label=(\alph*)]
\item Como $e^{i0X}=1$, temos $\varphi_X(0)=\mathbb{E}[1]=1$. 

\item Escrevendo $\varphi_X$ em partes real e imaginária, como na
\Cref{def:va-complexa}, e usando as propriedades de paridade das funções seno e cosseno temos
\[
\varphi_X(-t)=\mathbb{E}\bigl[\cos(-tX)\bigr]+i\,\mathbb{E}\bigl[\sin(-tX)\bigr]
=\mathbb{E}\bigl[\cos(tX)\bigr]-i\,\mathbb{E}\bigl[\sin(tX)\bigr]
=\overline{\varphi_X(t)}.
\]

\item Aplicando o \Cref{lem:esperanca-modulo}: $|\mathbb{E}[Z]|\leqslant\mathbb{E}\bigl[|Z|\bigr]$
com $Z=e^{itX}$ e lembrando que $|e^{itX}|=1$, obtemos
\[
|\varphi_X(t)|=\bigl|\mathbb{E}\bigl[e^{itX}\bigr]\bigr|
\leqslant\mathbb{E}\bigl[|e^{itX}|\bigr]=1,
\]
de modo que a imagem de $\varphi_X$ está contida no disco unitário fechado.

\item Pela linearidade da esperança e novamente pelo \Cref{lem:esperanca-modulo}, temos
\[
\begin{aligned}
\bigl|\varphi_X(t+h)-\varphi_X(t)\bigr|
&=\bigl|\mathbb{E}\bigl[e^{i(t+h)X}-e^{itX}\bigr]\bigr|
\leqslant\mathbb{E}\bigl[\bigl|e^{i(t+h)X}-e^{itX}\bigr|\bigr]
\\[0.3cm]
&=\mathbb{E}\bigl[\bigl|e^{itX}\bigl(e^{ihX}-1\bigr)\bigr|\bigr]
=\mathbb{E}\bigl[|e^{ihX}-1|\bigr],
\end{aligned}
\]
pois $|e^{itX}|=1$. Para mostrar que $t\longmapsto \varphi_{X}(t)$ é uniformemente 
contínua em toda reta, o ponto essencial é que o lado direito da desigualdade acima não depende de $t$.
Para mostrar que esta expressão tende a $0$ quando $h\to0$ consideramos a variável aleatória
$Y_h\equiv|e^{ihX}-1|=2\bigl|\sin(hX/2)\bigr|$. Note que esta variável aleatória
complexa é limitada por $2$ e tende a $0$ pontualmente, pela continuidade do seno e porque $hX(\omega)\to0$, para cada
$\omega$. Logo, segue do Teorema da Convergência Dominada que $\mathbb{E}\bigl[|e^{ihX}-1|\bigr]\to0$.

Como esse majorante independe de $t$, a mesma cota vale em toda a reta. Dado
$\varepsilon>0$, escolha $\delta>0$ tal que
$\mathbb{E}\bigl[|e^{ihX}-1|\bigr]<\varepsilon$ sempre que $|h|<\delta$; então
$|\varphi_X(t+h)-\varphi_X(t)|<\varepsilon$ para todo $t\in\mathbb{R}$ com
$|h|<\delta$. Isto é a continuidade uniforme de $\varphi_X$ em $\mathbb{R}$ e, em
particular, a continuidade em cada ponto.

\item Aqui basta usar as propriedades da exponencial complexa, a linearidade da
esperança e o fato de que $e^{itb}$ é constante, logo pode sair da esperança:
\[
\mathbb{E}\bigl[e^{it(aX+b)}\bigr]
=\mathbb{E}\bigl[e^{itb}\,e^{i(ta)X}\bigr]
=e^{itb}\,\mathbb{E}\bigl[e^{i(ta)X}\bigr]
=e^{itb}\,\varphi_X(at).
\]
Assim, se $Y\equiv aX+b$, a função característica de $Y$ é
$\varphi_Y(t)=e^{itb}\varphi_X(at)$. Essa é a regra de transformação afim da
função característica: uma translação multiplica $\varphi_X$ por uma exponencial
complexa de módulo $1$ e uma dilatação reescala o argumento. Vamos usá-la
repetidamente nos exemplos, para reduzir cada cálculo ao caso de uma variável
aleatória com parâmetros normalizados.
\end{enumerate}
\end{proof}

\subsection{Somas de Variáveis Aleatórias Independentes}\label{subsec:fc-soma}

A razão principal para introduzir funções características é o resultado a seguir:
ele converte somas de variáveis aleatórias independentes em produtos de funções.
Produtos são muito mais fáceis de analisar do que convoluções, de modo que,
sempre que a independência estiver presente, podemos trocar o problema analítico
de calcular uma convolução pelo problema algébrico de multiplicar funções
conhecidas.

\begin{theorem}\label{thm:fc-soma-independentes}
Sejam $X_1$ e $X_2$ variáveis aleatórias independentes, com funções
características
\[
\varphi_{X_j}(t) = \mathbb{E}\big[e^{itX_j}\big], \qquad t \in \mathbb{R}, \quad j = 1,2.
\]
Então $X_1 + X_2$ tem função característica $\varphi_{X_1}(t)\varphi_{X_2}(t)$, isto é,
\[
\mathbb{E}\big[e^{it(X_1+X_2)}\big] = \varphi_{X_1}(t)\,\varphi_{X_2}(t),
\qquad t \in \mathbb{R}.
\]
\end{theorem}

\begin{proof}
Fixemos $t \in \mathbb{R}$. O argumento se organiza em três etapas.

\smallskip

Pela propriedade básica da exponencial
complexa, $e^{a+b} = e^{a}e^{b}$ para quaisquer $a,b \in \mathbb{C}$. Tomando
$a = itX_1$ e $b = itX_2$, obtemos a identidade pontual
\[
e^{it(X_1+X_2)} = e^{itX_1}e^{itX_2},
\]
válida em todos os pontos do espaço de probabilidade.

\smallskip

As variáveis aleatórias $e^{itX_j}$, com $j=1,2$, são integráveis, já que $|e^{itX_j}| = 1$ e portanto
$\mathbb{E}[|e^{itX_j}|] = 1 < \infty$. Elas também são
independentes: cada uma é uma função mensurável de $X_j$, e funções mensuráveis
de variáveis aleatórias independentes, consideradas em pares, continuam
independentes. Sendo assim, vale a
regra do produto para esperanças de variáveis independentes e integráveis,
\[
\mathbb{E}\big[e^{itX_1}e^{itX_2}\big]
= \mathbb{E}\big[e^{itX_1}\big]\,\mathbb{E}\big[e^{itX_2}\big].
\]

\smallskip

Juntando as observações acima e usando as definições de $\varphi_{X_1}$ e $\varphi_{X_2}$,
temos
\[
\mathbb{E}\big[e^{it(X_1+X_2)}\big]
= \mathbb{E}\big[e^{itX_1}e^{itX_2}\big]
= \mathbb{E}\big[e^{itX_1}\big]\,\mathbb{E}\big[e^{itX_2}\big]
= \varphi_{X_1}(t)\,\varphi_{X_2}(t).
\]
Como $t \in \mathbb{R}$ era arbitrário a prova do teorema está completa.
\end{proof}

O resultado se estende sem esforço a qualquer número finito de parcelas.

\begin{corollary}\label{cor:fc-soma-n}
Se $X_1,\dots,X_n$ são variáveis aleatórias independentes, com funções
características $\varphi_{X_1},\dots,\varphi_{X_n}$, então $X_1+\cdots+X_n$ tem função
característica $\prod_{j=1}^{n}\varphi_{X_j}(t)$.
\end{corollary}

\begin{proof}
Indução em $n$: o caso $n = 1$ é trivial e, supondo o resultado válido para
$n-1$ parcelas, aplicamos o \Cref{thm:fc-soma-independentes} à variável
$X_1+\cdots+X_{n-1}$, que é independente de $X_n$ (por ser função mensurável de
$X_1,\dots,X_{n-1}$), e obtemos
\[
\mathbb{E}\big[e^{it(X_1+\cdots+X_n)}\big]
= \Bigl(\prod_{j=1}^{n-1}\varphi_{X_j}(t)\Bigr)\varphi_{X_n}(t)
= \prod_{j=1}^{n}\varphi_{X_j}(t).
\]
\end{proof}

\section{Exemplos}\label{sec:exemplos}

Vamos agora calcular a função característica das distribuições mais usadas no
curso. 

\subsection{Dois exemplos discretos}\label{subsec:exemplos-discretos}

Vamos começar por dois exemplos discretos bem simples: uma variável aleatória
com distribuição de Rademacher e uma variável aleatória com distribuição de
Poisson. 

\begin{example}[Distribuição de Rademacher]\label{ex:rademacher}
Seja $X$ uma variável aleatória com distribuição de Rademacher, isto é, suponha
que $X$ assuma apenas os valores $1$ e $-1$, com função de probabilidade dada por 
$\mathbb{P}(X=1)=\mathbb{P}(X=-1)=\tfrac12$. Como
$\operatorname{Im}X=\{1,-1\}$ é finito, a esperança da variável aleatória
complexa $e^{itX}$ é uma soma finita. Pela definição de função característica e
pela fórmula para a esperança de uma função de uma variável aleatória discreta, temos
$\tfrac12$:
\begin{align*}
\varphi_X(t)
&= \mathbb{E}\bigl[e^{itX}\bigr]
\\[0.3cm]
&= \sum_{x\in\operatorname{Im}X} e^{itx}\,\mathbb{P}(X=x)
\\[0.3cm]
&= \frac{e^{it}+e^{-it}}{2}.
\end{align*}
Pela fórmula de Euler, $e^{it}=\cos t+i\sin t$ e
$e^{-it}=\cos t-i\sin t$. Os senos se cancelam e os cossenos se somam, logo
\[
\varphi_X(t)=\cos t, \qquad t\in\mathbb{R}.
\]

É importante registrar três fatos que podem ser observados neste exemplo. 
Primeiro, a verificação da validade do item (a) do
\Cref{thm:propriedades-fc}: em $t=0$, temos $\varphi_X(0)=\cos 0=1$, como deve
ser. Segundo, $\varphi_X$ é real e par, o que reflete a simetria da distribuição:
trocar $X$ por $-X$ não altera a lei e corresponde a trocar $t$ por $-t$ dentro
da exponencial. Esse ponto será retomado no \Cref{cor:fc-real}, onde veremos que
a nulidade da parte imaginária da função característica é equivalente à simetria da distribuição.
Terceiro, $|\varphi_X(t)|=|\cos t|\leqslant 1$, o que verifica a validade da cota geral
$|\varphi_X(t)|\leqslant\varphi_X(0)$.
\end{example}

\begin{example}[Distribuição de Poisson]\label{ex:poisson}
Seja $X\sim\mathrm{Poisson}(\lambda)$, com $\lambda>0$, isto é,
\[
\mathbb{P}(X=k)=e^{-\lambda}\, \frac{\lambda^k}{k!}, \qquad k=0,1,2,\ldots.
\]

Como $X$ assume infinitos valores, a esperança que define $\varphi_X(t)$ é uma
série infinita. Antes de calculá-la, verificamos que ela converge absolutamente.
Essa convergência é suficiente para fatorar a constante $e^{-\lambda}$ e
identificar a soma com a série da exponencial sem nos preocuparmos com a ordem
dos termos.

Como $|e^{itk}|=1$ para todo $k$ e todo $t$, a série dos módulos é limitada por
\begin{align*}
\sum_{k=0}^{\infty}
\Bigl|e^{-\lambda}\,\frac{\lambda^{k}e^{itk}}{k!}\Bigr|
&= e^{-\lambda}\sum_{k=0}^{\infty}\frac{\lambda^{k}}{k!}
= e^{-\lambda}e^{\lambda}=1<\infty.
\end{align*}
Esta desigualdade mostra que a convergência da série é uniforme em $t$. 
Para o cálculo abaixo, porém, basta saber que a série acima converge absolutamente.

Com a questão da convergência absoluta resolvida, podemos manipular 
livremente a série que aparece abaixo. Já que
$e^{itk}=(e^{it})^{k}$ podemos acoplar o fator oscilante $e^{itk}$ ao peso
de Poisson $\lambda^{k}/k!$ e concluir que
\begin{align*}
\varphi_X(t)
&= \sum_{k=0}^{\infty} e^{itk}\,\mathbb{P}(X=k)
\\[0.3cm]
&= e^{-\lambda}\sum_{k=0}^{\infty}\frac{\bigl(\lambda e^{it}\bigr)^{k}}{k!}
\\[0.3cm]
&= e^{-\lambda}\exp\bigl(\lambda e^{it}\bigr)
\\[0.3cm]
&= \exp\bigl(\lambda(e^{it}-1)\bigr).
\end{align*}

Neste caso a verificação do item (a) do \Cref{thm:propriedades-fc} é imediata, já que em $t=0$
temos 
\[ 
\varphi_X(0)=\exp\bigl(\lambda(1-1)\bigr)=1.
\] 
Diferentemente do exemplo de
Rademacher, a distribuição de Poisson não é simétrica: ela só assume valores não
negativos e, para $\lambda>0$, não está concentrada em $\{0\}$. Por isso não há
razão para esperar que $\varphi_X$ seja real. Separando as partes real e
imaginária da exponencial complexa, obtemos
\[
\varphi_X(t)
=\exp\bigl(\lambda(\cos t-1)\bigr)
\Bigl(\cos(\lambda\sin t)+i\sin(\lambda\sin t)\Bigr),
\]
e portanto
$\operatorname{Im}\varphi_X(t)=\exp(\lambda(\cos t-1))\sin(\lambda\sin t)$,
que não se anula para todo $t$.
\end{example}

\subsection{A distribuição normal}\label{subsec:exemplo-normal}

\begin{example}[Distribuição normal padrão]\label{ex:normal}
Seja $Z$ uma variável aleatória com distribuição normal padrão, isto é, com
densidade $f(x)=\frac{1}{\sqrt{2\pi}}e^{-x^2/2}$, $x\in\mathbb{R}$. Vamos mostrar que a
função característica de $Z$ é dada por
\[
\varphi_Z(t)=\mathbb{E}\bigl[e^{itZ}\bigr]
=
\frac{1}{\sqrt{2\pi}}
\int_{-\infty}^{\infty}e^{itx}\,e^{-x^2/2}\,dx
=e^{-t^2/2}, \qquad \forall t\in\mathbb{R}.
\]

Apresentamos dois argumentos para esse cálculo. O primeiro é apenas
de natureza heurística ele é um bom dispositivo para identificar o valor correto de
uma certa integral. Entrenta, usando um pouco de Análise Complexa: Fórmula Integral de Cauchy
e deformação de contornos é possível torná-lo em uma prova da identidade. 
O segundo, argumento demonstra o resultado por um caminho completamente diferente relacionando o 
valor da integral a um problema de valor inicial.
\end{example}

\begin{proof}[Argumento heurístico]
Primeiro passo é reescrever o integrando em uma única exponencial,
$$e^{itx}e^{- \frac{x^2}{2}  }=e^{-(\frac{x^2}{2} -itx)}.$$ 
Em seguida, completamos o quadrado no expoente.
Como $(x-it)^2=x^2-2itx-t^2$, dividindo por $2$ e reorganizando os termos
obtemos a identidade
\[
\frac{x^2}{2}-itx=\frac{(x-it)^2}{2}+\frac{t^2}{2}.
\]
Substituindo essa expressão na integral e usando a propriedade
$e^{a+b}=e^ae^b$ da exponencial, segue que
\begin{align*}
\varphi_Z(t)
&=
\frac{1}{\sqrt{2\pi}}
\int_{-\infty}^{\infty}
\exp\biggl(-\frac{(x-it)^2}{2}-\frac{t^2}{2}\biggr)\,dx
\\[0.3cm]
&=e^{-t^2/2}\int_{-\infty}^{\infty}\frac{1}{\sqrt{2\pi}}
\exp\biggl(-\frac{(x-it)^2}{2}\biggr)\,dx ,
\end{align*}
pois o fator $e^{-t^2/2}$ não depende de $x$ e pode sair da integral.

A integral que resta é a integral sobre toda a reta da densidade de uma ``normal''
com ``média $it$'' e variância $1$. Se a média fosse um número real $a$, essa
integral valeria $1$, pela normalização da densidade da normal. 
\end{proof}

\begin{proof}[Prova formal]
Primeiro usar um argumento de paridade para eliminar a parte imaginária. Escrevendo
$e^{itx}=\cos(tx)+i\sin(tx)$ e usando que a função
$x\mapsto\sin(tx)\frac{1}{\sqrt{2\pi}}e^{-x^2/2}$ é ímpar e integrável em
$\mathbb{R}$ concluímos que 
\[
\varphi_Z(t)=\int_{-\infty}^{\infty}\cos(tx)\,\frac{1}{\sqrt{2\pi}}e^{-x^2/2}\,dx
=\mathbb{E}\bigl[\cos(tZ)\bigr].
\]
O integrando agora é real, e a majoração a seguir mostra que é legítimo trocar a ordem
da esperança e da derivação com respeito a $t$.

\bigskip 

De fato, note que para cada $t$ fixado, a derivada em $t$ do integrando é
$-x\sin(tx)\frac{1}{\sqrt{2\pi}}e^{-x^2/2}$, cujo módulo é dominado por
$|x|\frac{1}{\sqrt{2\pi}}e^{-x^2/2}$. Como
\[
\int_{-\infty}^{\infty}|x|\frac{1}{\sqrt{2\pi}}e^{-x^2/2}\,dx=\mathbb{E}\bigl[|Z|\bigr]<\infty,
\]
a hipótese (iii) está satisfeita com $Y=|Z|$, pois
$\bigl|\partial_t\cos(tZ)\bigr|=|Z|\,|\sin(tZ)|\leqslant|Z|$, e a majorante
não depende de $t$. Pelo \Cref{thm:derivacao-integral}, aplicado a
$g(t,z)=\cos(tz)$ e $X=Z$, podemos então derivar dentro da esperança e
obter
\[
\varphi_Z'(t)=\int_{-\infty}^{\infty}\bigl(-x\sin(tx)\bigr)
\frac{1}{\sqrt{2\pi}}e^{-x^2/2}\,dx .
\]
Para calcular essa integral, observe que
\[
\frac{d}{dx}\Bigl[\frac{1}{\sqrt{2\pi}}e^{-x^2/2}\Bigr]
=-x\,\frac{1}{\sqrt{2\pi}}e^{-x^2/2},
\]
de modo que $-x\,\frac{1}{\sqrt{2\pi}}e^{-x^2/2}$ é a derivada da densidade normal
padrão e o integrando tem a forma $\sin(tx)f'(x)$. Integrando por partes em um
intervalo $[-R,R]$ e fazendo $R\to\infty$,
\begin{align*}
\varphi_Z'(t)
&=\int_{-\infty}^{\infty}\sin(tx)\,\frac{d}{dx}
\Bigl[\frac{1}{\sqrt{2\pi}}e^{-x^2/2}\Bigr]\,dx
\\[0.3cm]
&=-\int_{-\infty}^{\infty}t\cos(tx)\,\frac{1}{\sqrt{2\pi}}e^{-x^2/2}\,dx
=-t\,\varphi_Z(t),
\end{align*}
porque o termo de fronteira se anula, já que $e^{-x^2/2}\to 0$, quando
$|x|\to\infty$ enquanto $|\sin(tx)|\leqslant 1$, e a última integral é
exatamente a expressão de $\varphi_Z(t)$ obtida acima.

Obtemos assim a equação diferencial linear de primeira ordem
$\varphi_Z'(t)=-t\varphi_Z(t)$, válida para todo $t\in\mathbb{R}$. Multiplicando
pelo fator integrante $e^{t^2/2}$,
\[
\frac{d}{dt}\Bigl[\varphi_Z(t)e^{t^2/2}\Bigr]
=e^{t^2/2}\bigl(\varphi_Z'(t)+t\varphi_Z(t)\bigr)=0 ,
\]
logo $\varphi_Z(t)e^{t^2/2}$ é constante em $\mathbb{R}$. Essa constante vale
$\varphi_Z(0)=1$, pois $f$ é uma densidade. Assim,
$\varphi_Z(t)e^{t^2/2}=1$ para todo $t$ e, portanto,
$\varphi_Z(t)=e^{-t^2/2}$, como queríamos demonstrar.
\end{proof}

\medskip
Combinando o resultado com o item (e) do \Cref{thm:propriedades-fc}, obtemos a
função característica de qualquer membro da família normal. Se $X$ tem
distribuição normal com média $\mu$ e variância $\sigma^2>0$, então
$X$ tem a mesma distribuição que $\mu+\sigma Z$, com $Z$ normal padrão, e o item
(e) aplicado com $a=\sigma$ e $b=\mu$ fornece
\begin{align*}
\varphi_X(t)
&=\mathbb{E}\bigl[e^{it(\mu+\sigma Z)}\bigr]
=e^{i\mu t}\,\mathbb{E}\bigl[e^{i(\sigma t)Z}\bigr]
\\[0.3cm]
&=e^{i\mu t}\,\varphi_Z(\sigma t)
=\exp\Bigl(i\mu t-\frac{\sigma^2t^2}{2}\Bigr), \qquad t\in\mathbb{R}.
\end{align*}
Uma distribuição normal com média $\mu$ e variância $\sigma^2$ tem, portanto,
função característica $\exp(i\mu t-\sigma^2t^2/2)$.

O mesmo tipo de escalonamento se aplica aos outros exemplos, de modo que
frequentemente daremos a função característica apenas para um membro
conveniente da família, deixando os demais parâmetros para uma conta como a
acima.

\subsection{Quatro exemplos com densidade}\label{subsec:exemplos-continuos}

Nos dois primeiros exemplos desta subseção, a densidade é nula fora de um
intervalo limitado; por isso, as integrais que definem a função característica
se reduzem a integrais sobre intervalos finitos. Nos outros dois, o suporte é
ilimitado, mas o decaimento exponencial da densidade garante a convergência das
integrais impróprias. Os cálculos usam integração direta, convolução e simetria.

\begin{example}[Distribuição uniforme]\label{ex:uniforme}
Seja $X$ uma variável aleatória com densidade
$f(x)=\frac{1}{b-a}\mathds{1}_{(a,b)}(x)$, com $-\infty<a<b<\infty$. Vamos
mostrar que
\[
\varphi_X(t)=\frac{1}{(b-a)}\frac{e^{itb}-e^{ita}}{it}, \qquad
t\in\mathbb{R}\setminus\{0\},
\]
e que $\varphi_X(0)=1$.
\end{example}

\begin{proof}
A primitiva de $x\mapsto e^{\lambda x}$ é $x\mapsto e^{\lambda x}/\lambda$
quando $\lambda\neq0$. No caso complexo, essa identidade segue separando
$e^{\lambda x}$ em partes real e imaginária e aplicando a regra do produto a
cada uma delas. Tomando $\lambda=it$ com $t\neq0$ e aplicando o teorema
fundamental do cálculo, obtemos
\[
\int_a^b e^{itx}\,dx
=\frac{e^{itb}-e^{ita}}{it}.
\]
Dividindo pela massa total $b-a$ da densidade, segue a expressão anunciada para
$\varphi_X(t)$.

Resta o valor em $t=0$. A fórmula acima tem uma indeterminação aparente nesse
ponto, mas a definição dá diretamente
\[
\varphi_X(0)=\int_{\mathbb{R}}f(x)\,dx=1.
\]
Além disso, quando $t\to0$, a expansão
$e^{itx}=1+itx+O(t^{2}x^{2})$ fornece
\[ e^{itb}-e^{ita}=it(b-a)+O(t^{2}),\] 
de modo que o quociente tende a $1$.
Portanto, a expressão obtida para $t\neq0$ se estende continuamente a $t=0$
com o valor $1$, como exige a continuidade de $\varphi_X$.
\end{proof}

No caso simétrico $a=-c$ e $b=c$, com $c>0$, o numerador é
$e^{itc}-e^{-itc}=2i\sin(ct)$ e o denominador é $2ict$. Assim,
\[
\varphi_X(t)=\frac{\sin(ct)}{ct}, \qquad t\neq0,
\]
com $\varphi_X(0)=1$. Para $c=1$, isso dá $\sin t/t$. Vale comparar com
\Cref{ex:rademacher}. Se $R$ tem distribuição de Rademacher, então
$\varphi_R(t)=\cos t$, pois a média de $e^{itx}$ é tomada apenas sobre os
pontos $-1$ e $1$. Aqui a média é tomada sobre todo o intervalo $(-c,c)$, o
que produz $\sin(ct)/(ct)$. Nesse sentido, a distribuição uniforme é a versão
contínua da distribuição de Rademacher.

\medskip

\begin{example}[Distribuição triangular]\label{ex:triangular}
Seja $X$ uma variável aleatória com densidade
$f(x)=(1-|x|)\mathds{1}_{(-1,1)}(x)$. Vamos mostrar que
\[
\varphi_X(t)=\frac{2\,(1-\cos t)}{t^{2}}, \qquad t\neq0,
\]
e que $\varphi_X(0)=1$. O nome vem do gráfico de $f$, que sobe de $0$ a $1$ e
desce.
\end{example}

\begin{proof}
A ideia é representar a densidade triangular como a densidade de uma soma de
duas variáveis uniformes independentes. Sejam $U$ e $V$ independentes, ambas
com densidade constante igual a $1$ em $(-1/2,1/2)$ e nula fora dele. Pela
fórmula da convolução para a soma de variáveis independentes,
\[
f_{U+V}(s)=\int_{-1/2}^{1/2}\mathds{1}_{(-1/2,1/2)}(s-y)\,dy,
\]
porque a densidade de $V$ vale $1$ em todo o intervalo de integração. O
integrando é $1$ exatamente quando $y\in(-1/2,1/2)$ e $s-y\in(-1/2,1/2)$, isto
é, quando $y\in(s-1/2,s+1/2)$. Portanto, a integral é o comprimento da
interseção $(-1/2,1/2)\cap(s-1/2,s+1/2)$. Para $s\in[0,1)$, essa interseção é
$(s-1/2,1/2)$, de comprimento $1-s$; para $s\in(-1,0]$, o cálculo simétrico dá
$1+s$; e, para $|s|\geqslant1$, ela é vazia. Logo
$f_{U+V}(s)=(1-|s|)\mathds{1}_{(-1,1)}(s)$, que é precisamente a densidade de
$X$. Como a função característica depende apenas da distribuição, para obter
$\varphi_X$ basta calcular $\varphi_{U+V}$.

Por \Cref{ex:uniforme}, com $a=-1/2$ e $b=1/2$,
\[
\varphi_U(t)=\frac{e^{it/2}-e^{-it/2}}{it}
=\frac{2\sin(t/2)}{t}, \qquad t\neq0,
\]
e $\varphi_U(0)=1$. Pela independência,
\Cref{thm:fc-soma-independentes} diz que a função característica da soma é o
produto das funções características. Assim,
\[
\varphi_{U+V}(t)
=\varphi_U(t)\varphi_V(t)
=\left(\frac{2\sin(t/2)}{t}\right)^{2}, \qquad t\neq0.
\]
Como $X$ e $U+V$ têm a mesma distribuição, $\varphi_X=\varphi_{U+V}$.

A identidade $\cos 2\theta=1-2\sin^{2}\theta$, com $\theta=t/2$, fornece
$\sin^{2}(t/2)=\bigl(1-\cos t\bigr)/2$. Substituindo no quadrado, obtemos
\[
\varphi_X(t)=\frac{4(1-\cos t)}{2t^{2}}
=\frac{2(1-\cos t)}{t^{2}}.
\]

Em $t=0$, a primeira expressão tende a $1$, pois $\sin(t/2)\sim t/2$, e a
segunda também, pois $1-\cos t\sim t^{2}/2$. Os dois limites coincidem com
$\varphi_X(0)=\int_{\mathbb{R}}f(x)\,dx=1$, de modo que as expressões obtidas
para $t\neq0$ se estendem continuamente à origem. Como a densidade é simétrica
em torno de $0$, era de se esperar que $\varphi_X$ fosse real e par.
\end{proof}

\medskip

\begin{example}[Distribuição exponencial]\label{ex:exponencial}
Seja $X$ uma variável aleatória com densidade
$f(x)=e^{-x}\mathds{1}_{(0,\infty)}(x)$. Vamos mostrar que
\[
\varphi_X(t)=\frac{1}{1-it}, \qquad t\in\mathbb{R}.
\]
\end{example}

\begin{proof}
Por definição,
\[
\varphi_X(t)=\int_{0}^{\infty}e^{itx}e^{-x}\,dx
=\int_{0}^{\infty}e^{(it-1)x}\,dx.
\]
A integral converge absolutamente, pois $|e^{itx}e^{-x}|=e^{-x}$ é integrável
em $(0,\infty)$. Podemos, portanto, calculá-la como limite de integrais em
$[0,R]$. Como $it-1\neq0$ para todo $t$ real, a primitiva é
$e^{(it-1)x}/(it-1)$, e
\[
\int_{0}^{R}e^{(it-1)x}\,dx=\frac{e^{(it-1)R}-1}{it-1}.
\]
O termo de fronteira superior se anula, já que
$\bigl|e^{(it-1)R}\bigr|=\bigl|e^{itR}\bigr|e^{-R}=e^{-R}\to0$ quando
$R\to\infty$: o fator oscilante tem módulo $1$ e o decaimento exponencial
domina. Passando ao limite,
\[
\varphi_X(t)=\frac{0-1}{it-1}=\frac{1}{1-it}.
\]
A fórmula vale para todo $t$ real, sem necessidade de extensão por continuidade,
porque $1-it$ nunca se anula. Em particular, $\varphi_X(0)=1$.
\end{proof}

\medskip

\begin{example}[Exponencial bilateral]\label{ex:exponencial-bilateral}
Seja $X$ uma variável aleatória com densidade $f(x)=\tfrac12e^{-|x|}$,
$x\in\mathbb{R}$. Sua distribuição é conhecida como \emph{exponencial
bilateral} (ou distribuição de Laplace), e sua função característica é
\[
  \varphi_X(t) \;=\; \frac{1}{1+t^2}, \qquad t\in\mathbb{R}.
\]
\end{example}

\begin{proof}
Como a expressão de $f$ muda em $x=0$, começamos separando a integral nesse
ponto. Para $t\in\mathbb{R}$, a definição de função característica dá
\[
  \varphi_X(t)
  = \frac12\int_{-\infty}^{0} e^{itx}e^{x}\,dx
    + \frac12\int_{0}^{\infty} e^{itx}e^{-x}\,dx.
\]
As integrais convergem absolutamente, pois $|e^{itx}|=1$ e $e^{-|x|}$ é
integrável em $\mathbb{R}$. Na primeira integral, fazemos a mudança de
variável $u=-x$. Assim, usando $e^{itu}+e^{-itu}=2\cos(tu)$, obtemos
\[
  \varphi_X(t)
  = \frac12\int_{0}^{\infty}\bigl(e^{-itu}+e^{itu}\bigr)e^{-u}\,du
  = \int_{0}^{\infty}e^{-u}\cos(tu)\,du.
\]
Para calcular a última integral, aproveitamos o cálculo do
\Cref{ex:exponencial}. Como $e^{-u}\cos(tu)$ é a parte real de
$e^{(it-1)u}$, temos
\[
  \varphi_X(t)
  = \operatorname{Re}\int_{0}^{\infty}e^{(it-1)u}\,du
  = \operatorname{Re}\frac{1}{1-it}
  = \operatorname{Re}\frac{1+it}{1+t^{2}}
  = \frac{1}{1+t^{2}}.
\]
O cancelamento das partes imaginárias reflete a simetria da densidade, e a
função característica encontrada é real e par. Em particular, $\varphi_X(0)=1$.
\end{proof}

\section{A Fórmula de Inversão}\label{sec:inversao}

Nos exemplos da seção anterior, partimos da distribuição de uma variável
aleatória e calculamos sua função característica. Vamos agora percorrer o
caminho inverso: dada a função característica de $X$, queremos recuperar
probabilidades associadas a $X$ e, com elas, sua função de distribuição.

\subsection{O Teorema de Inversão}\label{subsec:teorema-inversao}

A fórmula a seguir recupera a probabilidade de um intervalo, com uma correção
nos extremos quando eles têm probabilidade positiva.

\begin{theorem}[Fórmula de inversão]\label{thm:inversao}
Seja $X$ uma variável aleatória real com função característica $\varphi_X$.
Então, para quaisquer $a<b$,
\begin{equation}\label{eq:inversao}
\lim_{T\to\infty}\frac{1}{2\pi}
\int_{-T}^{T}\frac{e^{-ita}-e^{-itb}}{it}\,\varphi_X(t)\,dt
=
\mathbb{P}(a<X<b)
+\frac12\,\mathbb{P}(X=a)
+\frac12\,\mathbb{P}(X=b).
\end{equation}
\end{theorem}

O quociente no integrando tem uma descontinuidade apenas aparente em $t=0$:
seu limite nesse ponto é $b-a$, como veremos na prova. Assim, a integral sobre
$[-T,T]$ está bem definida. O teorema afirma, em particular, que o limite
dessas integrais existe.


\medskip 

Seja $F_X(x)=\mathbb{P}(X\leqslant x)$ a função de distribuição de $X$.
Para interpretar o resultado em termos de $F_X$, recorde que
$F_X(x^-)=\mathbb{P}(X<x)$ e que
$\mathbb{P}(X=x)=F_X(x)-F_X(x^-)$. Definindo
\[
F_X^*(x)\equiv\frac{F_X(x^-)+F_X(x)}{2},
\]
e usando $\mathbb{P}(a<X<b)=F_X(b^-)-F_X(a)$, podemos escrever o lado direito
de \eqref{eq:inversao} como $F_X^*(b)-F_X^*(a)$. Se $F_X$ é contínua em $a$
e $b$, as probabilidades dos dois extremos são nulas, e o lado direito de 
\eqref{eq:inversao} se reduz a
\[
\mathbb{P}(a<X\leqslant b)=F_X(b)-F_X(a).
\]

\begin{proof}
A ideia é escrever a integral em \eqref{eq:inversao} como a esperança de uma
função de $X$. Vamos calcular o limite dessa função para cada valor possível
de $X$ e, depois, passar o limite para dentro da esperança.

Antes disso, recordemos uma consequência da convergência quase certa que
usaremos duas vezes. Se $Y_n\to Y$ quase certamente e $|Y_n|\leqslant M$
para uma constante $M$, então $\mathbb{E}[Y_n]\to\mathbb{E}[Y]$. De fato,
$|Y|\leqslant M$ quase certamente e, para todo $\varepsilon>0$,
\[
\mathbb{E}\bigl[|Y_n-Y|\bigr]
\leqslant
\varepsilon+2M\,\mathbb{P}\bigl(|Y_n-Y|>\varepsilon\bigr).
\]
A probabilidade à direita tende a $0$, pois a convergência quase certa
implica convergência em probabilidade. Como $\varepsilon$ é arbitrário,
$\mathbb{E}[|Y_n-Y|]\to0$ e, pela desigualdade para o módulo da esperança,
$\mathbb{E}[Y_n]\to\mathbb{E}[Y]$. O mesmo argumento vale para variáveis
complexas limitadas. Aqui usamos as propriedades de linearidade e
monotonicidade da esperança para variáveis aleatórias limitadas.

Comecemos pelo quociente que aparece na fórmula. Para $t\neq0$,
\[
\frac{e^{-ita}-e^{-itb}}{it}=\int_a^b e^{-ity}\,dy.
\]
A expressão à direita vale $b-a$ em $t=0$ e é contínua em $t$. Definindo
$k(t)$ por essa integral, obtemos uma função contínua que coincide com o
quociente para $t\neq0$. Além disso, para todo $t\in\mathbb{R}$,
\[
|k(t)|\leqslant\int_a^b|e^{-ity}|\,dy=b-a.
\]

Para $T>0$, definimos
\[
K_T(x)\equiv\frac{1}{2\pi}\int_{-T}^{T}k(t)e^{itx}\,dt,
\qquad x\in\mathbb{R}.
\]
Como $\varphi_X(t)=\mathbb{E}[e^{itX}]$, queremos trocar a esperança com a
integral em $t$. Para justificar essa troca, aproximamos a integral por somas
de Riemann. Para cada valor de $X$, o integrando é contínuo em $t$, e as
somas são limitadas em módulo por $T(b-a)/\pi$. Pela propriedade de
convergência de variáveis limitadas que acabamos de recordar, podemos passar
o limite das somas para dentro da esperança. A linearidade da esperança e a
continuidade de $\varphi_X$ dão, portanto,
\begin{equation}\label{eq:IT}
\frac{1}{2\pi}\int_{-T}^{T}k(t)\varphi_X(t)\,dt
=\mathbb{E}\bigl[K_T(X)\bigr].
\end{equation}

Vamos agora calcular $K_T(x)$. Para $t\neq0$, temos
\[
k(t)e^{itx}
=\frac{e^{it(x-a)}-e^{it(x-b)}}{it}.
\]
Ao separar as partes real e imaginária, a parte que contém os cossenos é
ímpar em $t$ e se cancela na integral sobre $[-T,T]$. Resta
\[
K_T(x)
=\frac{1}{2\pi}\int_{-T}^{T}
\frac{\sin\bigl(t(x-a)\bigr)-\sin\bigl(t(x-b)\bigr)}{t}\,dt.
\]
Os quocientes se estendem continuamente a $t=0$.

Para usar a integral de Dirichlet, definimos
\[
S(u)\equiv\int_0^u\frac{\sin v}{v}\,dv,
\qquad u\in\mathbb{R},
\]
com $\sin v/v=1$ em $v=0$. O integrando é par, logo $S$ é ímpar. Pelo
\Cref{lem:dirichlet}, $S(u)\to\pi/2$ quando $u\to+\infty$. Já que esta função é ímpar,
$S(u)\to-\pi/2$ quando $u\to-\infty$. Além disso, $S$ é contínua e tem
limites finitos nos dois extremos da reta, de modo que existe uma constante
$M<\infty$, satisfazendo $|S(u)|\leqslant M$, para todo $u\in\mathbb{R}$.

Para $\theta\neq0$, o integrando $\sin(\theta t)/t$ é par em $t$. Assim,
fazendo a mudança de variável $v=\theta t$, com $t=v/\theta$ e
$dt=dv/\theta$, obtemos
\begin{equation}\label{eq:RT}
\begin{aligned}
\int_{-T}^{T}\frac{\sin(\theta t)}{t}\,dt
&=2\int_0^T\frac{\sin(\theta t)}{t}\,dt
\\[0.3cm]
&=2\int_0^{T\theta}\frac{\sin v}{v/\theta}\,\frac{dv}{\theta}
\\[0.3cm]
&=2\int_0^{T\theta}\frac{\sin v}{v}\,dv
=2S(T\theta).
\end{aligned}
\end{equation}

Observe que se $\theta<0$, então o limite superior $T\theta$ é negativo, e mesmo assim a mudança de
variável continua válida com os limites de integração na ordem indicada.
Para $\theta=0$, ambos os lados de \eqref{eq:RT} são nulos.

Aplicando \eqref{eq:RT} com $\theta=x-a$ e $\theta=x-b$, temos
\begin{align*}
K_T(x)
&=\frac{1}{2\pi}\left[
\int_{-T}^{T}\frac{\sin\bigl(t(x-a)\bigr)}{t}\,dt
-\int_{-T}^{T}\frac{\sin\bigl(t(x-b)\bigr)}{t}\,dt
\right]
\\[0.3cm]
&=\frac{1}{2\pi}\bigl[2S\bigl(T(x-a)\bigr)-2S\bigl(T(x-b)\bigr)\bigr]
\\[0.3cm]
&=\frac{1}{\pi}\bigl[S\bigl(T(x-a)\bigr)-S\bigl(T(x-b)\bigr)\bigr].
\end{align*}

Para calcular o limite, observamos os sinais de $x-a$ e $x-b$. Se $x$ está
fora de $[a,b]$, esses sinais são iguais; se $a<x<b$, são opostos. Nos
extremos, um dos termos contém $S(0)=0$. Usando os limites de $S$, obtemos
\begin{align*}
x<a:\quad
\lim_{T\to\infty}K_T(x)
&=\frac{1}{\pi}\left(-\frac{\pi}{2}-\left(-\frac{\pi}{2}\right)\right)=0,
\\[0.3cm]
x=a:\quad
\lim_{T\to\infty}K_T(a)
&=\frac{1}{\pi}\left(0-\left(-\frac{\pi}{2}\right)\right)=\frac12,
\\[0.3cm]
a<x<b:\quad
\lim_{T\to\infty}K_T(x)
&=\frac{1}{\pi}\left(\frac{\pi}{2}-\left(-\frac{\pi}{2}\right)\right)=1,
\\[0.3cm]
x=b:\quad
\lim_{T\to\infty}K_T(b)
&=\frac{1}{\pi}\left(\frac{\pi}{2}-0\right)=\frac12,
\\[0.3cm]
x>b:\quad
\lim_{T\to\infty}K_T(x)
&=\frac{1}{\pi}\left(\frac{\pi}{2}-\frac{\pi}{2}\right)=0.
\end{align*}

Além disso, a desigualdade triangular e a cota $|S(u)|\leqslant M$ dão
\begin{align*}
|K_T(x)|
&=\frac{1}{\pi}\left|S\bigl(T(x-a)\bigr)-S\bigl(T(x-b)\bigr)\right|
\\[0.3cm]
&\leqslant\frac{1}{\pi}\left(
\left|S\bigl(T(x-a)\bigr)\right|
+\left|S\bigl(T(x-b)\bigr)\right|\right)
\leqslant\frac{2M}{\pi},
\qquad x\in\mathbb{R},\ T>0.
\end{align*}

Seja $(T_n)$ uma sequência arbitrária de números positivos com $T_n\to\infty$.
Pelos limites calculados acima,
\[
K_{T_n}(X)\longrightarrow
\mathds{1}_{\{a<X<b\}}
+\tfrac12\mathds{1}_{\{X=a\}}
+\tfrac12\mathds{1}_{\{X=b\}}.
\]
As variáveis $K_{T_n}(X)$ são limitadas pela mesma constante $2M/\pi$.
Aplicando a convergência de variáveis limitadas e usando \eqref{eq:IT},
obtemos
\begin{align*}
\lim_{n\to\infty}\frac{1}{2\pi}
\int_{-T_n}^{T_n}k(t)\varphi_X(t)\,dt
&=\lim_{n\to\infty}\mathbb{E}\bigl[K_{T_n}(X)\bigr]
\\[0.3cm]
&=\mathbb{E}\left[
\mathds{1}_{\{a<X<b\}}
+\tfrac12\mathds{1}_{\{X=a\}}
+\tfrac12\mathds{1}_{\{X=b\}}
\right]
\\[0.3cm]
&=\mathbb{P}(a<X<b)
+\tfrac12\mathbb{P}(X=a)
+\tfrac12\mathbb{P}(X=b).
\end{align*}
Na última igualdade, usamos a linearidade da esperança e
$\mathbb{E}[\mathds{1}_A]=\mathbb{P}(A)$. Como a sequência $T_n\to\infty$
é arbitrária, o limite em \eqref{eq:inversao} existe e tem o valor anunciado.
\end{proof}

\begin{remark}[O truncamento simétrico]\label{rem:inversao-limite}
A integral em \eqref{eq:inversao} é tomada sobre $[-T,T]$. Para ver por que
essa escolha merece atenção, considere $X=0$ quase certamente, $a=0$ e $b=1$.
Então $\varphi_X(t)=1$. Já o termo que aparece na integral é da forma
\[
\frac{1-e^{-it}}{it}
=\frac{\sin t}{t}-i\,\frac{1-\cos t}{t}
\qquad \forall t\in \mathbb{R}\setminus \{0\}.
\]
A parte imaginária é ímpar e se cancela no intervalo simétrico. Pela integral
de Dirichlet, o lado esquerdo de \eqref{eq:inversao} tende a
\[
\frac{1}{2\pi}\,2\int_0^\infty\frac{\sin t}{t}\,dt=\frac12,
\]
que coincide com $\tfrac12\mathbb{P}(X=0)$.

Por outro lado, a integral da parte imaginária sobre $(0,T]$ não tem limite
finito. De fato,
\[
\int_1^T\frac{1-\cos t}{t}\,dt
=\log T-\int_1^T\frac{\cos t}{t}\,dt\longrightarrow\infty,
\]
pois a última integral converge, como se vê por integração por partes.
Assim, neste exemplo, o cancelamento entre os dois lados da origem é
necessário para o limite da fórmula.
\end{remark}

\medskip

A fórmula de inversão mostra que a função característica determina a função
de distribuição. De fato, suponha que $X$ e $Y$ tenham a mesma função
característica. Para quaisquer $a<b$, a fórmula dá
\[
F_X^*(b)-F_X^*(a)=F_Y^*(b)-F_Y^*(a).
\]
Quando $a\to-\infty$, tanto $F_X^*(a)$ quanto $F_Y^*(a)$ tendem a zero, pois
$0\leqslant F_X^*(a)\leqslant F_X(a)$, e o mesmo vale para $Y$. Logo
$F_X^*(b)=F_Y^*(b)$ para todo $b$.

Resta recuperar $F_X$ a partir de $F_X^*$. Se $y>x$, a monotonicidade de
$F_X$ dá
\[
F_X(x)\leqslant F_X(y^-)\leqslant F_X^*(y)\leqslant F_X(y).
\]
Fazendo $y\downarrow x$ e usando a continuidade à direita, obtemos
$F_X^*(y)\to F_X(x)$. O mesmo vale para $Y$, portanto $F_X=F_Y$. Concluímos
que duas variáveis aleatórias com a mesma função característica têm a mesma
distribuição.

\subsection{Duas Consequências Imediatas}\label{subsec:consequencias-inversao}

A fórmula de inversão tem duas consequências imediatas, que provamos a seguir:
uma propriedade de simetria em termos de $\varphi_X$ e a distribuição da soma
de variáveis normais centradas independentes.

\begin{corollary}[Função característica real]\label{cor:fc-real}
Seja $X$ uma variável aleatória real. Se $\varphi_X$ assume apenas valores
reais, então $X$ e $-X$ têm a mesma distribuição.
\end{corollary}

\begin{proof}
Calculemos a função característica de $-X$. Por definição,
\[
  \varphi_{-X}(t) = \mathbb{E}\bigl[e^{it(-X)}\bigr]
  = \mathbb{E}\bigl[e^{-itX}\bigr], \qquad t \in \mathbb{R}.
\]
Como $e^{-i\theta} = \overline{e^{i\theta}}$ para todo $\theta$ real,
a variável $e^{-itX}$ é o conjugado de $e^{itX}$. Tomando as esperanças das
partes real e imaginária, obtemos
\[
  \varphi_{-X}(t) = \overline{\mathbb{E}\bigl[e^{itX}\bigr]}
  = \overline{\varphi_X(t)},
\]
que é o item (b) de \Cref{thm:propriedades-fc}. Por hipótese
$\varphi_X(t)$ é real, logo vale
$\varphi_{-X}(t) = \varphi_X(t)$ para todo $t$; pela unicidade dada pelo
\Cref{thm:inversao}, duas variáveis com a mesma função característica têm
a mesma distribuição, de modo que $X \overset{d}{=} -X$.
\end{proof}

\begin{corollary}[Soma de normais centradas independentes]\label{cor:soma-normais}
Sejam $X_1$ e $X_2$ independentes, com $X_i$ normal de média $0$ e variância
$\sigma_i^2$, $i = 1, 2$. Então $X_1 + X_2$ é normal de média $0$ e variância
$\sigma_1^2 + \sigma_2^2$.
\end{corollary}

\begin{proof}
Pela independência, o \Cref{thm:fc-soma-independentes} dá
$\varphi_{X_1+X_2} = \varphi_{X_1}\varphi_{X_2}$, e o \Cref{ex:normal}
fornece $\varphi_{X_i}(t) = \exp(-\sigma_i^2 t^2/2)$. Logo, usando
$e^{a}e^{b} = e^{a+b}$,
\[
  \varphi_{X_1+X_2}(t)
  = \exp\Bigl(-\frac{\sigma_1^2 t^2}{2}\Bigr)
    \exp\Bigl(-\frac{\sigma_2^2 t^2}{2}\Bigr)
  = \exp\Bigl(-\frac{(\sigma_1^2+\sigma_2^2)\,t^2}{2}\Bigr).
\]
Pelo \Cref{ex:normal}, o membro direito é a função característica de uma
variável $Y$ normal de média $0$ e variância $\sigma_1^2+\sigma_2^2$; como
$\varphi_{X_1+X_2}=\varphi_Y$, a unicidade do \Cref{thm:inversao} 
garante que $X_1+X_2\overset{d}{=}Y$.
\end{proof}

\medskip

O argumento do \Cref{cor:soma-normais} é o protótipo do uso das funções
características: diretamente, seria preciso calcular a convolução
\[
  f_{X_1+X_2}(z) = \int_{-\infty}^{\infty}
  \frac{1}{\sqrt{2\pi\sigma_1^2}}
  \exp\Bigl(-\frac{(z-y)^2}{2\sigma_1^2}\Bigr)
  \frac{1}{\sqrt{2\pi\sigma_2^2}}
  \exp\Bigl(-\frac{y^2}{2\sigma_2^2}\Bigr)\,dy
\]
e completar o quadrado no expoente. O caminho pelas funções características
troca essa integral por uma multiplicação de exponenciais: a independência vira
produto e a unicidade da inversão converte identidades entre funções em
identidades entre distribuições. Esse é o mecanismo dos teoremas limites.

\subsection{O Caso Integrável}\label{subsec:caso-integravel}

A fórmula de inversão toma uma forma mais simples quando
$\int_{\mathbb{R}}|\varphi_X(t)|\,dt<\infty$. Nesse caso, veremos que $X$
possui uma densidade, que também pode ser recuperada a partir de $\varphi_X$.
As integrais sobre toda a reta nesta subseção são integrais impróprias do
cálculo.

\begin{theorem}[Densidade quando $\varphi_X$ é integrável]\label{thm:densidade-fc-integravel}
Seja $X$ uma variável aleatória com função característica $\varphi_X$. Se
\[
\int_{-\infty}^{\infty}|\varphi_X(t)|\,dt<\infty,
\]
então $X$ possui uma densidade contínua e limitada $f$, dada por
\[
f(y)=\frac{1}{2\pi}\int_{-\infty}^{\infty}
e^{-ity}\varphi_X(t)\,dt,
\qquad y\in\mathbb{R}.
\]
Em particular, para todos $a<b$,
\[
\mathbb{P}(a<X\leqslant b)=F_X(b)-F_X(a)=\int_a^b f(y)\,dy.
\]
\end{theorem}

\begin{proof}
A ideia é usar a integrabilidade de $\varphi_X$ para substituir o limite no
\Cref{thm:inversao} por uma integral sobre toda a reta. Primeiro mostraremos
que $X$ não tem probabilidade positiva em nenhum ponto. Em seguida,
escreveremos a probabilidade de um intervalo como a integral da função $f$
do enunciado.

Como vimos na prova da fórmula de inversão, vale a seguinte desigualdade
\begin{equation}\label{eq:cota-nucleo-inversao}
\left|\frac{e^{-ita}-e^{-itb}}{it}\right|
=\left|\int_a^b e^{-ity}\,dy\right|
\leqslant b-a,
\qquad a<b.
\end{equation}
Além do mais, a função que aparece dentro do valor absoluto 
a esquerda da desigualdade acima pode ser estendida por
continuidade ao valor $b-a$, quando $t=0$. 

O módulo do integrando em
\eqref{eq:inversao} é, portanto, limitado por
$(b-a)|\varphi_X(t)|$. Como essa função é integrável, a integral imprópria
sobre toda a reta converge absolutamente. A fórmula de inversão fica
\begin{equation}\label{eq:inversao-integravel}
\mathbb{P}(a<X<b)
+\frac12\,\mathbb{P}(X=a)
+\frac12\,\mathbb{P}(X=b)
=
\frac{1}{2\pi}\int_{-\infty}^{\infty}
\frac{e^{-ita}-e^{-itb}}{it}\,\varphi_X(t)\,dt.
\end{equation}
Pela mesma cota, se
\[
C\equiv\frac{1}{2\pi}\int_{-\infty}^{\infty}|\varphi_X(t)|\,dt,
\]
então, para quaisquer $a<b$,
\[
\frac12\,\mathbb{P}(X=a)
\leqslant
\mathbb{P}(a<X<b)
+\frac12\,\mathbb{P}(X=a)
+\frac12\,\mathbb{P}(X=b)
\leqslant C(b-a).
\]
Fixado $a$, fazemos $b\downarrow a$ e concluímos que
$\mathbb{P}(X=a)=0$. Como $a$ é arbitrário, $F_X$ é contínua em toda a reta.

Definimos agora a função proposta no enunciado por
\[
f(y)\equiv\frac{1}{2\pi}\int_{-\infty}^{\infty}
e^{-ity}\varphi_X(t)\,dt.
\]
A integral converge absolutamente e $|f(y)|\leqslant C$ para todo $y$.
Para verificar a continuidade de $f$, considere primeiro as integrais
truncadas
\[
f_T(y)\equiv\frac{1}{2\pi}\int_{-T}^{T}e^{-ity}\varphi_X(t)\,dt.
\]
Para cada $T$, a função $f_T$ é contínua, pois o integrando é contínuo em
$(t,y)$ e o intervalo de integração é limitado. Além disso,
\[
|f(y)-f_T(y)|
\leqslant\frac{1}{2\pi}
\left(\int_{-\infty}^{-T}|\varphi_X(t)|\,dt
+\int_T^{\infty}|\varphi_X(t)|\,dt\right),
\]
e o lado direito tende a zero quando $T\to\infty$, independentemente de $y$.
Logo $f_T\to f$ uniformemente, e $f$ é contínua. A função $f$ também é
real: pela identidade $\varphi_X(-t)=\overline{\varphi_X(t)}$, as
contribuições de $t$ e $-t$ em cada integral truncada são conjugadas.

Vamos relacionar $f$ com as probabilidades de $X$. Como o integrando é
contínuo no retângulo limitado $[-T,T]\times[a,b]$, podemos trocar a ordem
das duas integrais usuais e escrever
\begin{align*}
\int_a^b f_T(y)\,dy
&=\frac{1}{2\pi}\int_{-T}^{T}
\left(\int_a^b e^{-ity}\,dy\right)\varphi_X(t)\,dt\\[0.3cm]
&=\frac{1}{2\pi}\int_{-T}^{T}
\frac{e^{-ita}-e^{-itb}}{it}\,\varphi_X(t)\,dt.
\end{align*}
Quando $T\to\infty$, a convergência uniforme de $f_T$ assegura que
$\int_a^b f_T(y)\,dy\to\int_a^b f(y)\,dy$. Pelo
\Cref{thm:inversao} e pela ausência de átomos, o lado
direito converge para $\mathbb{P}(a<X\leqslant b)$. Portanto,
\[
\mathbb{P}(a<X\leqslant b)=\int_a^b f(y)\,dy.
\]

Como $f$ é contínua e sua integral sobre qualquer intervalo é não negativa,
temos $f(y)\geqslant0$ para todo $y$. De fato, se $f$ fosse negativa em um
ponto, pela continuidade ela seria negativa em um intervalo ao redor desse
ponto, e sua integral sobre o intervalo também seria negativa. 

Por fim, para
$a<x$,
\[
F_X(x)-F_X(a)=\int_a^x f(y)\,dy.
\]
Fazendo $a\to-\infty$ e usando $F_X(a)\to0$, obtemos
$F_X(x)=\int_{-\infty}^x f(y)\,dy$. Quando $x\to\infty$, essa integral tende
a $1$, pois $F_X(x)\to1$. Logo $f$ é uma densidade de $X$.
\end{proof}

Se já conhecemos uma densidade contínua $g$ de $X$, então ela coincide em
todos os pontos com a densidade $f$ fornecida pelo teorema. De fato, $f$ e
$g$ têm a mesma integral sobre cada intervalo; como ambas são contínuas,
isso só é possível quando $f=g$ em toda a reta. Usaremos essa observação nos
exemplos seguintes.

A hipótese de integrabilidade de $\varphi_X$ é suficiente para a existência de
densidade, mas não é necessária. No \Cref{ex:exponencial}, por exemplo, $X$
tem densidade e
\[
|\varphi_X(t)|=\frac{1}{\sqrt{1+t^2}},
\]
cuja integral sobre $\mathbb{R}$ diverge por comparação com $1/|t|$ quando
$|t|$ é grande.

\subsection{Duas funções características obtidas por inversão}\label{subsec:exemplos-inversao}

Aplicando a fórmula de inversão do \Cref{thm:densidade-fc-integravel} às funções
características obtidas no \Cref{ex:triangular} e no \Cref{ex:exponencial-bilateral},
obtemos duas novas funções características. Nos dois casos o integrando da fórmula de
inversão é, a menos de uma constante, uma função característica já calculada na
\Cref{sec:exemplos}, e a identidade, lida na direção oposta, diz que a densidade
recuperada tem por função característica a densidade de partida.

\medskip

\begin{example}[Distribuição de Pólya]\label{ex:distribuicao-polya}
A função $f(x)=(1-\cos x)/(\pi x^{2})$, para $x\neq0$, com $f(0)=1/(2\pi)$, é uma
densidade de probabilidade em $\mathbb{R}$. Se $X$ tem densidade $f$, sua função
característica é
\[
\varphi_X(t)=\bigl(1-|t|\bigr)^{+}, \qquad t\in\mathbb{R},
\]
onde $a^{+}\equiv\max\{a,0\}$ denota a parte positiva de $a$.
\end{example}

\begin{proof}
Seja $Z$ uma variável aleatória com a distribuição triangular do
\Cref{ex:triangular}. Sua função característica é
$\varphi_Z(s)=2(1-\cos s)/s^{2}$ para $s\neq0$, com
$\varphi_Z(0)=1$. Ela é integrável: perto da origem é contínua com limite
$1$, pois $1-\cos s\sim s^{2}/2$, e para $|s|\geqslant1$ vale
$\bigl|\varphi_Z(s)\bigr|\leqslant4/s^{2}$. As hipóteses do
\Cref{thm:densidade-fc-integravel} estão satisfeitas. Como a densidade $f_Z$ é
contínua, a observação após o teorema dá, para todo $y\in\mathbb{R}$,
\[
f_Z(y)
=\frac{1}{2\pi}\int_{\mathbb{R}}\frac{2\,(1-\cos s)}{s^{2}}\,e^{-isy}\,ds
=\int_{\mathbb{R}}\frac{1-\cos s}{\pi s^{2}}\,e^{-isy}\,ds
=\bigl(1-|y|\bigr)^{+},
\]
em que a segunda igualdade absorve a constante e a última é o valor da densidade dado
no \Cref{ex:triangular}.

Fazendo $y=0$, segue que
$\int_{\mathbb{R}}(1-\cos s)/(\pi s^{2})\,ds=1$, pois o integrando é limitado perto de
$0$ e majorado por $2/(\pi s^{2})$ longe de $0$. A função $f$ do enunciado, que
coincide com esse integrando fora da origem, é portanto não negativa e tem integral
$1$; o valor $f(0)=1/(2\pi)$ apenas a torna contínua e não altera a integral.

Trocando agora os nomes das variáveis, com $s=x$ para a variável de integração e
$y=-t$ para o ponto de avaliação, e usando que $\bigl(1-|y|\bigr)^{+}$ é par,
\[
\int_{\mathbb{R}}\frac{1-\cos x}{\pi x^{2}}\,e^{ixt}\,dx
=\bigl(1-|t|\bigr)^{+}, \qquad t\in\mathbb{R}.
\]
O lado esquerdo é a função característica de $X$ avaliada em $t$, pela própria
definição, logo $\varphi_X(t)=\bigl(1-|t|\bigr)^{+}$ e, em particular, $\varphi_X(0)=1$.
Note que $\bigl(1-|t|\bigr)^{+}$ é contínua, vale $1$ em $t=0$ e se anula para
$|t|\geqslant1$: seu gráfico é um triângulo de base $[-1,1]$ e altura $1$, o mesmo
perfil da densidade do \Cref{ex:triangular} com o papel das variáveis trocado.
\end{proof}

A distribuição deste exemplo não tem nome consagrado na literatura. O nome que usamos
se deve ao papel que essa densidade desempenha em um critério de construção de funções
características devido a Pólya.

\medskip

\begin{example}[Distribuição de Cauchy]\label{ex:cauchy}
A função $f(x)=1/\bigl(\pi(1+x^{2})\bigr)$, $x\in\mathbb{R}$, é uma densidade de
probabilidade. Se $X$ tem densidade $f$, sua função característica é
\[
\varphi_X(t)=e^{-|t|}, \qquad t\in\mathbb{R}.
\]
\end{example}

\begin{proof}
Seja $Z$ uma variável aleatória com a distribuição exponencial bilateral do
\Cref{ex:exponencial-bilateral}. Sua função característica é
$\varphi_Z(s)=1/(1+s^{2})$, que é integrável: é contínua e limitada em $[-1,1]$ e, para $|s|\geqslant1$, vale
$\varphi_Z(s)\leqslant1/s^{2}$. Como a densidade da exponencial bilateral é contínua, o
\Cref{thm:densidade-fc-integravel} fornece sua expressão em cada ponto. Multiplicando
a identidade por $2$, obtemos
\begin{align*}
\frac{1}{2\pi}\int_{\mathbb{R}}\frac{1}{1+s^{2}}\,e^{-isy}\,ds
=\tfrac12\,e^{-|y|}
\quad \Longrightarrow \quad 
\int_{\mathbb{R}}\frac{1}{\pi(1+s^{2})}\,e^{-isy}\,ds
=e^{-|y|}, \qquad \forall y\in\mathbb{R}.
\end{align*}
Fazendo $y=0$ na segunda identidade e lembrando que a primitiva de $1/(1+s^{2})$ é
$\arctan s$, o lado esquerdo vale
$\frac{1}{\pi}\bigl(\frac{\pi}{2}+\frac{\pi}{2}\bigr)=1$, cálculo de integração já
conhecido; logo o integrando é uma densidade de probabilidade. Trocando $s=x$ e
$y=-t$, e usando de novo que $e^{-|y|}$ é par,
\[
\int_{\mathbb{R}}\frac{1}{\pi(1+x^{2})}\,e^{ixt}\,dx=e^{-|t|},
\qquad t\in\mathbb{R},
\]
o que diz que $\varphi_X(t)=e^{-|t|}$ é a função característica de $X$;
em particular $\varphi_X(0)=1$.
\end{proof}

A distribuição de Cauchy é o exemplo clássico de distribuição sem média:
$\mathbb{E}[|X|]=\frac{2}{\pi}\int_{0}^{\infty}x/(1+x^{2})\,dx=\infty$,
pois a primitiva do integrando é $\tfrac12\log(1+x^{2})$, que diverge. Isso se reflete
em $t=0$: a função $\varphi_X(t)=e^{-|t|}$ é derivável em todo $t\neq0$, com
$\varphi_X'(t)=-\operatorname{sgn}(t)\,e^{-|t|}$, mas não é derivável em $t=0$, onde as
derivadas laterais valem $-1$, à direita, e $+1$, à esquerda. Como $\varphi_X>0$, o
módulo também não é derivável na origem: $|\varphi_X(t)|=e^{-|t|}$ tem o mesmo bico.

Em ambos os exemplos as densidades são simétricas e as funções características
resultantes são reais e pares, como previsto no \Cref{cor:fc-real}. Além disso,
nenhuma das duas funções características é derivável em $t=0$, o que é coerente com a ausência
de primeiro momento finito nessas distribuições; voltaremos a essa relação na
\Cref{sec:momentos}.

\section{O Teorema da Continuidade de~Lévy}\label{sec:continuidade}

Nas notas de aula sobre modos de convergência, vimos que $(X_n)_{n\in\mathbb{N}}$ converge em
distribuição para $X$, e escrevemos $X_n\Rightarrow X$, quando
\[
F_n(x)=\mathbb{P}(X_n\leqslant x)
\longrightarrow
F(x)=\mathbb{P}(X\leqslant x)
\]
em todo ponto $x$ em que $F$ é contínua. As variáveis não precisam estar
definidas no mesmo espaço de probabilidade: essa convergência depende apenas
de suas funções de distribuição.

Vamos relacionar a convergência em distribuição com a convergência das
funções características. Para o sentido mais delicado, precisaremos controlar
a probabilidade de as variáveis assumirem valores muito distantes da origem.

\begin{definition}[Sequência condensada]\label{def:sequencia-condensada}
Dizemos que uma sequência de variáveis aleatórias $(X_n)_{n\geqslant1}$ é
\emph{condensada} se, para todo $\varepsilon>0$, existe $M>0$ tal que
\[
\mathbb{P}(|X_n|>M)<\varepsilon
\qquad\text{para todo }n\geqslant1.
\]
\end{definition}

Em outras palavras, dado $\varepsilon>0$, podemos escolher um único intervalo
$[-M,M]$ que contenha $X_n$ com probabilidade maior que $1-\varepsilon$ para
todos os índices $n$. Esse tipo de sequência é conhecido também na literatura
como ``\emph{tight}''. Alguns autores traduzem essa expressão para o português
como ``sequência apertada'' ou ``sequência justa''.

\begin{theorem}[Teorema da Continuidade de Lévy]\label{thm:continuidade}
Sejam $X_1,X_2,\ldots$ variáveis aleatórias reais, e denote por
$\varphi_{X_n}$ a função característica de $X_n$, para cada $n$.

\begin{enumerate}[label=\textup{(\roman*)}]
\item Se $X_n\Rightarrow X$, então
$\varphi_{X_n}(t)\to\varphi_X(t)$ para todo $t\in\mathbb{R}$.

\item Se $\varphi_{X_n}(t)\to\varphi(t)$ para todo $t\in\mathbb{R}$, para uma
função $\varphi$ contínua em $0$, então existe uma variável aleatória $X$ tal que
$X_n\Rightarrow X$ e $\varphi=\varphi_X$. Além disso, $(X_n)$ é condensada.
\end{enumerate}
\end{theorem}

No item (ii), $\varphi$ é inicialmente apenas uma função limite. Uma das conclusões
do teorema é que esta função limite é função característica de alguma variável aleatória.

Na prova, usaremos dois fatos sobre convergência em distribuição. O primeiro
permite passar das funções de distribuição às esperanças de funções limitadas
e contínuas. O segundo mostra que uma sequência condensada possui uma
subsequência convergente em distribuição.

\begin{lemma}[Esperanças e convergência em distribuição]\label{lem:dist-esperancas-limitadas}
Se $X_n\Rightarrow X$ e $h:\mathbb{R}\to\mathbb{C}$ é contínua e limitada,
então
\[
\mathbb{E}[h(X_n)]\longrightarrow\mathbb{E}[h(X)].
\]
\end{lemma}

\begin{proof}
A ideia é aproximar $h$ por uma função constante em cada intervalo de uma
partição, como ilustra a \Cref{fig:aprox-funcao-escada-esperanca}. 

\begin{figure}[H]
\centering
\begin{tikzpicture}[
  x=1.2cm, y=1.2cm,
  >=Stealth,
  font=\small,
  declare function={
    h(\x) = 0.85 + 0.9*cos(deg(0.55*\x)) - 0.4*sin(deg(0.75*\x));
  }
]
  % Parâmetros fundamentais
  \def\R{3.5}
  \def\B{2.35}
  \def\xone{-2.1}
  \def\xtwo{-0.7}
  \def\xthree{0.7}
  \def\xfour{2.1}

  % Faixas sombreadas para as caudas fora de (-R, R]
  \fill[black!4, rounded corners=2pt] (-5.35, -2.45) rectangle (-\R, 2.55);
  \fill[black!4, rounded corners=2pt] (\R, -2.45) rectangle (5.35, 2.55);

  % Região sombreada de erro no intervalo destacado (x_{j-1}, x_j]
  \fill[teal!14] (\xtwo, {h(\xthree)})
    -- plot[domain=\xtwo:\xthree, samples=40] (\x, {h(\x)})
    -- cycle;

  % Cotas horizontais y = B e y = -B
  \draw[dashed, muted!70, line width=0.5pt] (-5.35, \B) -- (5.35, \B)
    node[right=2pt, font=\scriptsize, text=muted] {$y = B$};
  \draw[dashed, muted!70, line width=0.5pt] (-5.35, -\B) -- (5.35, -\B)
    node[right=2pt, font=\scriptsize, text=muted] {$y = -B$};

  % Eixos coordenados
  \draw[->, black!65, line width=0.6pt] (-5.4, 0) -- (5.6, 0)
    node[right=3pt] {$x$};
  \draw[->, black!65, line width=0.6pt] (0, -2.5) -- (0, 2.7)
    node[above=3pt] {$y$};

  % Marcas no eixo y para B e -B
  \draw[black!65, line width=0.55pt] (-0.08, \B) -- (0.08, \B);
  \node[left=3pt, font=\scriptsize, text=muted] at (0, 2.5) {$B$};
  \draw[black!65, line width=0.55pt] (-0.08, -\B) -- (0.08, -\B);
  \node[left=3pt, font=\scriptsize, text=muted] at (0, -2.5) {$-B$};

  % Curva contínua h(x)
  \draw[black!85, line width=1.15pt, domain=-5.2:5.2, samples=120, smooth]
    plot (\x, {h(\x)});
  \node[above right=1pt, font=\small, text=black!90]
    at (4.7, {h(4.7)}) {$y = h(x)$};

  % Linhas guias verticais tracejadas dos nós
  \draw[teal!45, densely dashed, line width=0.5pt]
    (-\R, 0) -- (-\R, {h(\xone)});
  \draw[teal!45, densely dashed, line width=0.5pt]
    (\xone, 0) -- (\xone, {h(\xtwo)});
  \draw[teal!45, densely dashed, line width=0.5pt]
    (\xtwo, 0) -- (\xtwo, {h(\xtwo)});
  \draw[teal!45, densely dashed, line width=0.5pt]
    (\xthree, 0) -- (\xthree, {h(\xthree)});
  \draw[teal!45, densely dashed, line width=0.5pt]
    (\xfour, 0) -- (\xfour, {h(\xfour)});
  \draw[teal!45, densely dashed, line width=0.5pt]
    (\R, 0) -- (\R, {h(\R)});

  % Função escada g(x):
  % 1) Cauda esquerda: x <= -R ==> g(x) = 0
  \draw[teal, line width=1.5pt] (-5.2, 0) -- (-\R, 0);
  \fill[teal] (-\R, 0) circle (2.2pt);

  % 2) Intervalos da partição (-R, R]:
  % Degrau 1: (-R, x_1], g = h(x_1)
  \draw[teal, line width=1.5pt] (-\R, {h(\xone)}) -- (\xone, {h(\xone)});
  \draw[teal, fill=white, line width=0.9pt] (-\R, {h(\xone)}) circle (2.2pt);
  \fill[teal] (\xone, {h(\xone)}) circle (2.2pt);

  % Degrau 2: (x_1, x_{j-1}], g = h(x_{j-1})
  \draw[teal, line width=1.5pt] (\xone, {h(\xtwo)}) -- (\xtwo, {h(\xtwo)});
  \draw[teal, fill=white, line width=0.9pt] (\xone, {h(\xtwo)}) circle (2.2pt);
  \fill[teal] (\xtwo, {h(\xtwo)}) circle (2.2pt);

  % Degrau 3: (x_{j-1}, x_j], g = h(x_j) [degrau destacado]
  \draw[teal, line width=1.5pt] (\xtwo, {h(\xthree)}) -- (\xthree, {h(\xthree)});
  \draw[teal, fill=white, line width=0.9pt] (\xtwo, {h(\xthree)}) circle (2.2pt);
  \fill[teal] (\xthree, {h(\xthree)}) circle (2.2pt);

  % Degrau 4: (x_j, x_{m-1}], g = h(x_{m-1})
  \draw[teal, line width=1.5pt] (\xthree, {h(\xfour)}) -- (\xfour, {h(\xfour)});
  \draw[teal, fill=white, line width=0.9pt] (\xthree, {h(\xfour)}) circle (2.2pt);
  \fill[teal] (\xfour, {h(\xfour)}) circle (2.2pt);

  % Degrau 5: (x_{m-1}, R], g = h(R)
  \draw[teal, line width=1.5pt] (\xfour, {h(\R)}) -- (\R, {h(\R)});
  \draw[teal, fill=white, line width=0.9pt] (\xfour, {h(\R)}) circle (2.2pt);
  \fill[teal] (\R, {h(\R)}) circle (2.2pt);

  % 3) Cauda direita: x > R ==> g(x) = 0
  \draw[teal, line width=1.5pt] (\R, 0) -- (5.2, 0);
  \draw[teal, fill=white, line width=0.9pt] (\R, 0) circle (2.2pt);

  % Rótulo de g(x)
  \node[above=2pt, font=\small, text=teal] at (-2.8, {h(\xone)}) {$y = g(x)$};

  % Ticks e nós no eixo x
  \foreach \x/\rot in {-\R/{-R=x_0}, \xone/{x_1}, \xtwo/{x_{j-1}}, \xthree/{x_j}, \xfour/{x_{m-1}}, \R/{x_m=R}} {
    \draw[black!70, line width=0.55pt] (\x, -0.08) -- (\x, 0.08);
    \node[below=4pt, font=\scriptsize, fill=white, inner sep=1pt] at (\x, 0) {$\rot$};
  }

  % Destaque do erro no intervalo (x_{j-1}, x_j]: |h(x) - g(x)| < varepsilon
  \draw[{Stealth[length=3.5pt]}-{Stealth[length=3.5pt]}, teal!90!black, line width=0.8pt]
    (-0.25, {h(-0.25)}) -- (-0.25, {h(\xthree)});
  \node[above=2pt, font=\scriptsize, text=teal!90!black]
    at (1.25, {h(-2.0)}) {$|h(x)-g(x)| < \varepsilon$};

  % Anotações nas regiões de cauda
  \node[font=\scriptsize, text=muted, align=center] at (-4.42, 1.3)
    {Fora de $(-R, R]$:\\[1pt]
     $g(x) = 0$\\[1pt]
     $|h(x)|\leqslant B$};
  \node[font=\scriptsize, text=muted, align=center] at (-4.42, -1.45)
    {Cauda esquerda:\\[1pt]
     $\mathbb{P}(X \leqslant -R) < \varepsilon$};

  \node[font=\scriptsize, text=muted, align=center] at (4.42, 1.3)
    {Fora de $(-R, R]$:\\[1pt]
     $g(x) = 0$\\[1pt]
     $|h(x)|\leqslant B$};
  \node[font=\scriptsize, text=muted, align=center] at (4.42, -1.45)
    {Cauda direita:\\[1pt]
     $\mathbb{P}(X > R) < \varepsilon$};

  % Chave indicando a partição de [-R, R]
  \draw[line width=0.55pt, decorate, decoration={brace, mirror, amplitude=5pt}]
    (-\R, -0.72) -- (\R, -0.72)
    node[midway, below=7pt, font=\small, fill=white, inner sep=2pt]
      {Partição de $[-R, R]$};

\end{tikzpicture}
\caption{Aproximação da função contínua e limitada $h$ pela função escada $g$ na partição de $[-R, R]$. Em cada degrau $(x_{j-1}, x_j]$, a função $g$ vale a constante $h(x_j)$, com extremidade esquerda aberta e direita fechada, o que garante $|h(x)-g(x)|<\varepsilon$. Fora de $(-R, R]$, a função $g$ se anula e o erro pontual é majorado por $B$. 
A probabilidade acumulada nas caudas é estritamente menor do que $\varepsilon$.}
\label{fig:aprox-funcao-escada-esperanca}
\end{figure}

Nesses intervalos, as probabilidades podem ser escritas por meio
das funções de distribuição.

Sejam $F_n$ e $F$ as funções de distribuição de $X_n$ e $X$, e seja $B$
uma constante com $|h(x)|\leqslant B$ para todo $x$. Dado $\varepsilon>0$,
escolhemos $R>0$ de modo que $-R$ e $R$ sejam pontos de continuidade de $F$ e
\[
\mathbb{P}(|X|\geqslant R)<\varepsilon.
\]
Essa escolha é possível porque as caudas de $X$ tendem a zero e $F$ tem no
máximo uma quantidade enumerável de pontos de descontinuidade. De fato, cada
descontinuidade é um salto de tamanho $\mathbb{P}(X=x)>0$, e para cada inteiro
$m\geqslant1$ há apenas uma quantidade finita de saltos de tamanho pelo menos
$1/m$.

Como $h$ é uniformemente contínua em $[-R,R]$, podemos escolher uma partição
\[
-R=x_0<x_1<\cdots<x_m=R
\]
cujos pontos sejam de continuidade de $F$ e tal que
$|h(x)-h(y)|<\varepsilon$ sempre que $x$ e $y$ pertencem ao mesmo intervalo
$(x_{j-1},x_j]$. Definimos
\[
g(x)\equiv\sum_{j=1}^{m}h(x_j)\mathds{1}_{(x_{j-1},x_j]}(x).
\]
Fora de $(-R,R]$, temos $g(x)=0$. Assim,
\begin{align*}
\bigl|\mathbb{E}[h(X_n)]-\mathbb{E}[g(X_n)]\bigr|
&\leqslant\mathbb{E}\bigl[|h(X_n)-g(X_n)|\bigr]\\[0.3cm]
&\leqslant\varepsilon+B\,\mathbb{P}\bigl(X_n\notin(-R,R]\bigr).
\end{align*}
A mesma estimativa vale para $X$. Além disso,
\begin{align*}
\mathbb{E}[g(X_n)]
&=\sum_{j=1}^{m}h(x_j)\,
\mathbb{P}(x_{j-1}<X_n\leqslant x_j)\\[0.3cm]
&=\sum_{j=1}^{m}h(x_j)\bigl(F_n(x_j)-F_n(x_{j-1})\bigr)\\[0.3cm]
&\longrightarrow
\sum_{j=1}^{m}h(x_j)\bigl(F(x_j)-F(x_{j-1})\bigr)
=\mathbb{E}[g(X)].
\end{align*}
A convergência vale porque todos os extremos da partição são pontos de
continuidade de $F$.

Por fim, a probabilidade de ficar fora de $(-R,R]$ também converge:
\[
\mathbb{P}\bigl(X_n\notin(-R,R]\bigr)
=F_n(-R)+1-F_n(R)
\longrightarrow
F(-R)+1-F(R)
\leqslant\mathbb{P}(|X|\geqslant R)<\varepsilon.
\]
Combinando as estimativas,
\[
\limsup_{n\to\infty}
\bigl|\mathbb{E}[h(X_n)]-\mathbb{E}[h(X)]\bigr|
\leqslant2\varepsilon+2B\varepsilon.
\]
Como $\varepsilon>0$ é arbitrário, a prova está completa.
\end{proof}

\begin{lemma}[Seleção de uma subsequência]\label{lem:selecao-condensada}
Toda sequência condensada $(X_n)_{n\in\mathbb{N}}$ possui uma subsequência que converge em
distribuição para alguma variável aleatória real.
\end{lemma}

\begin{proof}
Vamos construir primeiro uma candidata à função de distribuição limite.
Escreva $F_n(x)=\mathbb{P}(X_n\leqslant x)$ e enumere os racionais como
$q_1,q_2,\ldots$. Como $0\leqslant F_n(q_1)\leqslant1$, podemos escolher uma
subsequência ao longo da qual $F_n(q_1)$ converge. Dessa subsequência,
escolhemos outra ao longo da qual $F_n(q_2)$ converge, e assim por diante.

Tomando o $k$-ésimo termo da $k$-ésima subsequência, obtemos uma subsequência
$(X_{n_k})$ tal que, para cada $q\in\mathbb{Q}$, existe o limite
\[
G(q)\equiv\lim_{k\to\infty}F_{n_k}(q).
\]
A função $G$ é não decrescente nos racionais. Para estendê-la à reta,
definimos
\[
F(x)\equiv\inf_{\substack{q\in\mathbb{Q}\\q>x}}G(q),
\qquad x\in\mathbb{R}.
\]
A função $F$ é não decrescente e contínua à direita. Para verificar a
continuidade, dado $\varepsilon>0$, escolha $q>x$ racional com
$G(q)<F(x)+\varepsilon$. Se $x<y<q$, então
\[
F(x)\leqslant F(y)\leqslant G(q)<F(x)+\varepsilon.
\]

A condensação garante que $F$ tem os limites necessários em $-\infty$ e
$+\infty$. Dado $\varepsilon>0$, escolha $M$ com
$\mathbb{P}(|X_n|>M)<\varepsilon$ para todo $n$. Se $q<-M$ é racional,
$F_n(q)\leqslant\mathbb{P}(|X_n|>M)<\varepsilon$, logo $G(q)\leqslant\varepsilon$.
Se $q>M$ é racional, $F_n(q)\geqslant1-\varepsilon$, logo
$G(q)\geqslant1-\varepsilon$. Se $x<-M$, podemos escolher um racional $q$
com $x<q<-M$ e, portanto, $F(x)\leqslant G(q)\leqslant\varepsilon$. Se
$x>M$, todos os racionais $q>x$ satisfazem $G(q)\geqslant1-\varepsilon$, de
modo que $F(x)\geqslant1-\varepsilon$. Como $\varepsilon$ é arbitrário,
concluímos que
\[
\lim_{x\to-\infty}F(x)=0
\qquad\text{e}\qquad
\lim_{x\to\infty}F(x)=1.
\]
Pela caracterização das funções de distribuição apresentada nas notas de aula
sobre esse assunto, existe uma variável aleatória $X$ com função de
distribuição $F$.

Resta verificar que $X_{n_k}\Rightarrow X$. Seja $x$ um ponto de continuidade
de $F$ e tomemos $r<p<x<q$, com $p$ e $q$ racionais. Pela monotonicidade,
\[
F(r)\leqslant G(p)
\leqslant\liminf_{k\to\infty}F_{n_k}(x)
\leqslant\limsup_{k\to\infty}F_{n_k}(x)
\leqslant G(q)\leqslant F(q).
\]
As desigualdades das extremidades seguem da definição de $F$: como $p>r$,
temos $F(r)\leqslant G(p)$, e, como $G$ é não decrescente,
$G(q)\leqslant F(q)$. Fazendo $r\uparrow x$ e $q\downarrow x$, concluímos,
pela continuidade de $F$ em $x$, que $F_{n_k}(x)\to F(x)$. Portanto
$X_{n_k}\Rightarrow X$.
\end{proof}

\begin{proof}[Prova do \Cref{thm:continuidade}]
Para o item (i), fixamos $t\in\mathbb{R}$ e aplicamos o
\Cref{lem:dist-esperancas-limitadas} à função contínua e limitada
$h(x)=e^{itx}$. Obtemos
\[
\varphi_{X_n}(t)=\mathbb{E}[e^{itX_n}]
\longrightarrow
\mathbb{E}[e^{itX}]=\varphi_X(t).
\]

Passemos ao item (ii). A prova tem duas partes: primeiro mostraremos que
$(X_n)$ é condensada; depois identificaremos o limite de suas subsequências
convergentes.

\smallskip
\noindent\textbf{Uma cota para as caudas.}
Começamos com uma variável aleatória real $Y$ qualquer, de função
característica $\varphi_Y$. A ideia é usar os valores de $\varphi_Y$ perto de $0$ para
controlar a probabilidade de $|Y|$ ser grande.

Para $u>0$ e $x\neq0$, a paridade do cosseno dá
\begin{align*}
\frac{1}{u}\int_{-u}^{u}\bigl(1-\cos(tx)\bigr)\,dt
&=2-\frac{1}{u}\left.\frac{\sin(tx)}{x}\right|_{-u}^{u}\\[0.3cm]
&=2\left(1-\frac{\sin(ux)}{ux}\right).
\end{align*}
A identidade também vale em $x=0$, tomando $\sin(ux)/(ux)=1$ nesse ponto.
Substituindo $x$ por $Y$ e passando a esperança pela integral, obtemos
\begin{equation}\label{eq:continuidade-cauda-identidade}
\frac{1}{u}\int_{-u}^{u}
\bigl(1-\operatorname{Re}\varphi_Y(t)\bigr)\,dt
=2\mathbb{E}\left[1-\frac{\sin(uY)}{uY}\right].
\end{equation}
A troca é justificada por somas de Riemann, como na prova do
\Cref{thm:inversao}: para cada valor de $Y$, o integrando é contínuo em $t$,
e todas as somas são limitadas em módulo por uma constante. A convergência
de variáveis limitadas permite então passar o limite para dentro da esperança.

Para comparar o lado direito com uma probabilidade, usamos que
$|\sin v|\leqslant|v|$ para todo $v$. De fato, se $v\geqslant0$,
\[
|\sin v|=\left|\int_0^v\cos s\,ds\right|
\leqslant\int_0^v|\cos s|\,ds\leqslant v,
\]
e, para $v<0$, a desigualdade segue da imparidade do seno. Logo
$1-\sin v/v\geqslant0$, com o valor $0$ em $v=0$. Além disso, se
$|v|\geqslant2$, então
\[
\frac{\sin v}{v}\leqslant\frac{1}{|v|}\leqslant\frac12,
\qquad\text{e, portanto,}\qquad
2\left(1-\frac{\sin v}{v}\right)\geqslant1.
\]
Tomando $v=uY$, concluímos que
\[
\mathds{1}_{\{|Y|\geqslant2/u\}}
\leqslant2\left(1-\frac{\sin(uY)}{uY}\right).
\]
Tomando esperanças e usando \eqref{eq:continuidade-cauda-identidade}, obtemos
a cota
\begin{equation}\label{eq:cota-concentracao}
\mathbb{P}(|Y|\geqslant2/u)
\leqslant
\frac{1}{u}\int_{-u}^{u}
\bigl(1-\operatorname{Re}\varphi_Y(t)\bigr)\,dt.
\end{equation}

\smallskip
\noindent\textbf{A sequência é condensada.}
Por hipótese, $\varphi_{X_n}(t)\to\varphi(t)$ para todo $t$. Como
$\varphi_{X_n}(0)=1$, temos $\varphi(0)=1$. Fixado $\varepsilon>0$, a
continuidade de $\varphi$ em $0$ permite escolher $u>0$ tal que
\[
0\leqslant1-\operatorname{Re}\varphi(t)<\frac{\varepsilon}{4},
\qquad |t|\leqslant u.
\]
A desigualdade à esquerda segue de $|\varphi(t)|\leqslant1$, obtida passando
ao limite em $|\varphi_{X_n}(t)|\leqslant1$.

Para usar essa informação na cota \eqref{eq:cota-concentracao}, precisamos
passar ao limite nas integrais das funções $\varphi_{X_n}$. Considere uma
variável auxiliar $U$, uniforme em $[-u,u]$, e defina
\[
W_n=1-\operatorname{Re}\varphi_{X_n}(U),
\qquad
W=1-\operatorname{Re}\varphi(U).
\]
Não precisamos supor qualquer relação entre $U$ e as variáveis $X_n$, pois
as funções $\varphi_{X_n}$ são determinísticas. Como cada $\varphi_{X_n}$ é contínua,
$W_n$ é uma variável aleatória; seu limite pontual $W$ também é uma variável
aleatória. A convergência pontual de $\varphi_{X_n}$ dá $W_n\to W$; além disso,
$0\leqslant W_n\leqslant2$ e $0\leqslant W\leqslant\varepsilon/4$. Como $U$
tem densidade $1/(2u)$ em $[-u,u]$, a fórmula da esperança e a convergência
de variáveis limitadas fornecem
\[
\frac{1}{u}\int_{-u}^{u}
\bigl(1-\operatorname{Re}\varphi_{X_n}(t)\bigr)\,dt
=2\mathbb{E}[W_n]
\longrightarrow2\mathbb{E}[W]\leqslant\frac{\varepsilon}{2}.
\]
Logo existe $N$ tal que, para todo $n\geqslant N$,
\[
\mathbb{P}(|X_n|\geqslant2/u)
\leqslant
\frac{1}{u}\int_{-u}^{u}
\bigl(1-\operatorname{Re}\varphi_{X_n}(t)\bigr)\,dt
<\varepsilon.
\]
Para cada um dos índices $n<N$, podemos aumentar o limite $2/u$ até que a
probabilidade de $|X_n|$ ultrapassá-lo seja menor que $\varepsilon$. Como há
apenas um número finito desses índices, um mesmo $M$ serve para todos. Assim,
\[
\mathbb{P}(|X_n|>M)<\varepsilon
\qquad\text{para todo }n,
\]
e $(X_n)$ é condensada.

\smallskip
\noindent\textbf{Identificação do limite.}
Pelo \Cref{lem:selecao-condensada}, existe uma subsequência
$X_{n_k}\Rightarrow X$ para alguma variável aleatória $X$. Pelo item (i),
\[
\varphi_{X_{n_k}}(t)\longrightarrow\varphi_X(t)
\qquad\text{para todo }t\in\mathbb{R}.
\]
Mas, pela hipótese do item (ii), o mesmo limite é $\varphi(t)$. Logo
$\varphi=\varphi_X$.

Toda subsequência de $(X_n)$ continua sendo condensada, portanto possui uma
subsequência convergente em distribuição. O argumento anterior mostra que
qualquer limite assim obtido tem função característica $\varphi_X$. Pela
unicidade do \Cref{thm:inversao}, todos esses limites têm a mesma função de
distribuição $F_X$.

Resta mostrar que a sequência inteira converge. Se isso não ocorresse,
existiriam um ponto de continuidade $x$ de $F_X$, um número $\eta>0$ e uma
subsequência $(X_{n_j})$ com
\[
\bigl|\mathbb{P}(X_{n_j}\leqslant x)-F_X(x)\bigr|\geqslant\eta
\qquad\text{para todo }j.
\]
Essa subsequência possui, porém, uma subsequência convergente em distribuição
para uma variável com função de distribuição $F_X$. No ponto de continuidade
$x$, suas funções de distribuição devem convergir para $F_X(x)$, o que
contradiz a desigualdade. Portanto $X_n\Rightarrow X$.
\end{proof}

\begin{remark}[Por que a continuidade em $0$ é necessária]\label{rem:continuidade-origem}
Seja $Z$ uma variável normal padrão e considere $X_n=\sqrt{n}\,Z$. Então
$X_n$ tem distribuição normal de média $0$ e variância $n$, com função
característica
\[
\varphi_{X_n}(t)=e^{-nt^2/2},\qquad t\in\mathbb{R}.
\]
Para $t\neq0$, temos $e^{-nt^2/2}\to0$, enquanto $\varphi_{X_n}(0)=1$ para todo
$n$. O limite pontual é
\[
\varphi(t)=
\begin{cases}
1, & t=0,\\[0.2cm]
0, & t\neq0,
\end{cases}
\]
que não é contínuo na origem e, portanto, não pode ser uma função
característica.

Nesse exemplo, a falta de continuidade acompanha a dispersão das variáveis.
Para qualquer $M>0$,
\[
\mathbb{P}(|X_n|\leqslant M)
=\mathbb{P}\left(|Z|\leqslant\frac{M}{\sqrt n}\right)
\longrightarrow\mathbb{P}(Z=0)=0.
\]
Assim, $(X_n)$ não é condensada. Além disso, não pode convergir em
distribuição: pelo item (i) do \Cref{thm:continuidade}, a função característica
de um limite teria de coincidir com $\varphi$, mas toda função característica
é contínua em $0$.
\end{remark}

\section{Momentos e Derivadas}\label{sec:momentos}

Seja $X$ uma variável aleatória real. Nas notas de aula sobre momentos e funções
geradoras de momentos, vimos que as esperanças das potências de $X$ podem ser
recuperadas a partir das derivadas de
$M_X$ na origem, quando a função geradora existe em uma vizinhança de $0$.
Vamos investigar uma relação semelhante para a função característica
\[
\varphi_X(t)=\mathbb{E}[e^{itX}].
\]

A função característica está definida para todo $t\in\mathbb{R}$, pois
$|e^{itX}|=1$. Isso não garante, por si só, que ela seja derivável. Por
exemplo, no \Cref{ex:cauchy} encontramos uma variável sem primeiro momento
finito cuja função característica é $e^{-|t|}$, que não é derivável em
$t=0$.

Para entender o papel dos momentos, observamos que, ao derivar a exponencial
em relação a $t$, aparece um fator $X$:
\[
\frac{d}{dt}e^{itX}=iX e^{itX},
\qquad
\left|\frac{d}{dt}e^{itX}\right|=|X|.
\]
Assim, para justificar a passagem da derivada para dentro da esperança, uma
hipótese natural é que $\mathbb{E}[|X|]<\infty$. O resultado abaixo mostra
como o mesmo raciocínio se estende às derivadas de ordem superior.

\subsection{Derivadas, momentos e o decaimento das caudas}\label{subsec:momentos-derivadas}

Na prova do \Cref{thm:continuidade}, obtivemos uma cota que relaciona a função
característica às probabilidades de valores grandes de $|X|$: para $u>0$,
\[
\mathbb{P}(|X|\geqslant2/u)
\leqslant
\frac{1}{u}\int_{-u}^{u}
\bigl(1-\operatorname{Re}\varphi_X(t)\bigr)\,dt.
\]
Essa é a cota \eqref{eq:cota-concentracao}, escrita para a variável $X$. Ela
mostra que, quando a parte real de $\varphi_X$ está próxima de $1$ perto da
origem, a probabilidade de $|X|$ ser grande é pequena.

Os momentos também controlam essas probabilidades. De fato, se
$\mathbb{E}[|X|^n]<\infty$ para algum inteiro $n\geqslant1$, então, para
$r>0$,
\[
\mathds{1}_{\{|X|\geqslant r\}}\leqslant\frac{|X|^n}{r^n}.
\]
Tomando esperanças, obtemos
\[
\mathbb{P}(|X|\geqslant r)
\leqslant\frac{\mathbb{E}[|X|^n]}{r^n}.
\]
Vamos agora mostrar outra relação entre momentos e função característica:
quando o momento absoluto de ordem $n$ é finito, podemos derivar
$\varphi_X$ até a ordem $n$ e calcular suas derivadas por meio de esperanças.

\begin{theorem}[Derivadas da função característica]
\label{thm:derivadas-momentos}
Seja $X$ uma variável aleatória com função característica $\varphi_X$.
Se $\mathbb{E}[|X|^n]<\infty$ para algum inteiro $n\geqslant1$, então, para
cada inteiro $k$ com $1\leqslant k\leqslant n$, a derivada $\varphi_X^{(k)}$
existe, é contínua em $\mathbb{R}$ e satisfaz
\[
\varphi_X^{(k)}(t)
=\mathbb{E}\bigl[(iX)^k e^{itX}\bigr],
\qquad t\in\mathbb{R}.
\]
\end{theorem}

\begin{proof}
Vamos derivar a esperança uma vez de cada vez. Para aplicar o
\Cref{thm:derivacao-integral}, precisamos verificar que as variáveis obtidas
e as derivadas dos integrandos são integráveis.

Primeiro, a hipótese de momento absoluto de ordem $n$ garante os momentos
absolutos de ordens menores. Com efeito, se $1\leqslant j\leqslant n$, então
\[
|x|^j\leqslant1+|x|^n,\qquad x\in\mathbb{R}.
\]
Para $|x|\leqslant1$, temos $|x|^j\leqslant1$; para $|x|>1$, temos
$|x|^j\leqslant|x|^n$. Logo
\[
\mathbb{E}[|X|^j]\leqslant1+\mathbb{E}[|X|^n]<\infty,
\qquad 1\leqslant j\leqslant n.
\]

Comecemos pela primeira derivada. Para cada $t$, a variável $e^{itX}$ é
integrável, pois tem módulo $1$. Além disso,
\[
\frac{\partial}{\partial t}e^{itX}=iX e^{itX},
\qquad
\bigl|iX e^{itX}\bigr|=|X|.
\]
Como $\mathbb{E}[|X|]<\infty$, o \Cref{thm:derivacao-integral} permite
trocar a derivada com a esperança:
\[
\varphi_X'(t)
=\frac{d}{dt}\mathbb{E}[e^{itX}]
=\mathbb{E}[iX e^{itX}],
\qquad t\in\mathbb{R}.
\]

Suponha agora que, para algum $k<n$, já tenhamos obtido
\[
\varphi_X^{(k)}(t)=\mathbb{E}\bigl[(iX)^k e^{itX}\bigr].
\]
Para derivar mais uma vez, tomamos $g(t,x)=(ix)^k e^{itx}$. Então
\[
\frac{\partial g}{\partial t}(t,X)
=(iX)^{k+1}e^{itX},
\qquad
\left|\frac{\partial g}{\partial t}(t,X)\right|
=|X|^{k+1}.
\]
A variável $g(t,X)$ é integrável, pois seu módulo é $|X|^k$, e a derivada é
majorada pela variável integrável $|X|^{k+1}$. Aplicando novamente o
\Cref{thm:derivacao-integral}, concluímos que
\[
\varphi_X^{(k+1)}(t)
=\frac{d}{dt}\mathbb{E}\bigl[(iX)^k e^{itX}\bigr]
=\mathbb{E}\bigl[(iX)^{k+1}e^{itX}\bigr].
\]
A indução fornece, portanto, a fórmula para todas as ordens de $1$ a $n$.

Resta verificar a continuidade. Fixe $k\in\{1,\ldots,n\}$ e tome uma
sequência $t_m\to t$. Para cada resultado do experimento,
\[
(iX)^k e^{it_mX}\longrightarrow(iX)^k e^{itX},
\qquad
\bigl|(iX)^k e^{it_mX}\bigr|=|X|^k.
\]
Como $\mathbb{E}[|X|^k]<\infty$, a convergência dominada, aplicada às partes
real e imaginária, dá
\[
\varphi_X^{(k)}(t_m)
=\mathbb{E}\bigl[(iX)^k e^{it_mX}\bigr]
\longrightarrow
\mathbb{E}\bigl[(iX)^k e^{itX}\bigr]
=\varphi_X^{(k)}(t).
\]
Logo $\varphi_X^{(k)}$ é contínua para cada $k\leqslant n$.
\end{proof}

Tomando $t=0$ no teorema e usando $e^0=1$, obtemos
\[
\varphi_X^{(k)}(0)=i^k\,\mathbb{E}[X^k],
\qquad 1\leqslant k\leqslant n.
\]
Em particular, se $\mathbb{E}[|X|]<\infty$, então
\[
\varphi_X'(0)=i\,\mathbb{E}[X].
\]
Se $\mathbb{E}[X^2]<\infty$, temos também
\[
\varphi_X''(0)=i^2\,\mathbb{E}[X^2]=-\mathbb{E}[X^2].
\]
Nesse caso, a média é $\mathbb{E}[X]=-i\varphi_X'(0)$, e a variância é
\[
\begin{aligned}
\operatorname{Var}(X)
&=\mathbb{E}[X^2]-\bigl(\mathbb{E}[X]\bigr)^2
\\[0.3cm]
&=-\varphi_X''(0)-\bigl(-i\varphi_X'(0)\bigr)^2
=-\varphi_X''(0)+\bigl(\varphi_X'(0)\bigr)^2.
\end{aligned}
\]

Por exemplo, no \Cref{ex:rademacher} vimos que uma variável $R$ que assume
$-1$ e $1$ com probabilidades iguais tem função característica
$\varphi_R(t)=\cos t$. Como $\varphi_R'(0)=0$ e
$\varphi_R''(0)=-1$, as fórmulas dão
\[
\mathbb{E}[R]=0,
\qquad
\mathbb{E}[R^2]=1,
\]
em acordo com o cálculo direto a partir da distribuição de $R$.

A existência das derivadas também permite aproximar $\varphi_X$ perto da origem.
Se $\mathbb{E}[|X|^n]<\infty$, aplicamos a fórmula de Taylor às partes real e
imaginária de $\varphi_X$ e substituímos as derivadas pelos momentos:
\[
\varphi_X(t)
=1+\sum_{k=1}^{n}\frac{(it)^k}{k!}\,\mathbb{E}[X^k]
+o(t^n),
\qquad t\to0.
\]
Aqui $o(t^n)$ representa um resto cujo quociente por $t^n$ tende a zero
quando $t\to0$. Em particular, se $\mathbb{E}[X^2]<\infty$,
\[
\varphi_X(t)
=1+it\,\mathbb{E}[X]
-\frac{t^2}{2}\,\mathbb{E}[X^2]+o(t^2).
\]
Para essa conclusão, não precisamos supor que o terceiro momento seja finito.

Na próxima subseção, obteremos uma cota explícita para o resto da expansão da
exponencial. Ela permitirá estimar o erro antes de tomar a esperança e será
útil no estudo de somas de variáveis aleatórias.

\subsection{Uma cota para o resto de Taylor}\label{subsec:resto-taylor}

Em \Cref{thm:derivadas-momentos} vimos que as derivadas de $\varphi_X$ na
origem são os momentos de $X$. Para transformar essa informação em uma expansão
de Taylor, porém, não basta conhecer as derivadas: é preciso saber quão grande
é o erro cometido ao truncar a série. A cota a seguir resolve essa questão; ela
é uma afirmação puramente analítica, sobre a exponencial complexa, e por isso
vem antes de qualquer hipótese sobre $X$.

\begin{lemma}[Cota para o resto de Taylor]\label{lem:resto-taylor}
Para todo inteiro $n\geqslant0$ e todo $x\in\mathbb{R}$,
\[
\Bigl|e^{ix}-\sum_{m=0}^{n}\frac{(ix)^m}{m!}\Bigr|\leqslant\min\Bigl(\frac{|x|^{n+1}}{(n+1)!},\,\frac{2|x|^{n}}{n!}\Bigr).
\]
Quando $n=0$ a soma se reduz a $1$ e o segundo membro do mínimo deve ser lido
como $2$, de modo que a cota diz simplesmente $|e^{ix}-1|\leqslant\min(|x|,2)$.
\end{lemma}

Vale entender, antes da prova, por que o enunciado traz um mínimo de duas cotas.
O primeiro membro, $|x|^{n+1}/(n+1)!$, é a ordem de grandeza usual do resto de
Taylor de $e^{ix}$: ele descreve o erro como se descreveria o de uma série de
potências comum e é a cota adequada quando $|x|$ é pequeno. O segundo membro,
$2|x|^{n}/n!$, é pior por um fator da ordem de $|x|$ nessa região, mas cresce
como $|x|^{n}$ e não como $|x|^{n+1}$, sendo portanto o adequado para $|x|$
grande. É essa segunda cota, e não a primeira, que permitirá provar o teorema
central do limite sem supor o terceiro momento finito: ela exige apenas o
controle de $n$ momentos. Como as duas estimativas valem simultaneamente, vale
a menor delas.

\begin{proof}
A prova tem três etapas: obter uma identidade integral exata para o resto,
cotá-la de duas maneiras, cada uma eficiente em uma região, e tomar o mínimo.

\medskip
\textbf{Primeira etapa: a identidade do resto.}
Integrando por partes com $u=e^{is}$ e $dv=(x-s)^{n}ds$, de modo que
$du=ie^{is}ds$ e $v=-(x-s)^{n+1}/(n+1)$, obtemos
\[
\int_{0}^{x}(x-s)^{n}e^{is}ds=\Bigl[-\frac{(x-s)^{n+1}}{n+1}e^{is}\Bigr]_{s=0}^{s=x}+\frac{i}{n+1}\int_{0}^{x}(x-s)^{n+1}e^{is}ds .
\]
O termo de fronteira se anula em $s=x$, porque $(x-x)^{n+1}=0$, e vale
$-x^{n+1}/(n+1)$ em $s=0$; logo o colchete é $x^{n+1}/(n+1)$ e
\[
\int_{0}^{x}(x-s)^{n}e^{is}ds=\frac{x^{n+1}}{n+1}+\frac{i}{n+1}\int_{0}^{x}(x-s)^{n+1}e^{is}ds .
\]
A conta vale também para $x<0$: a integral é orientada e o teorema fundamental
do cálculo continua valendo nos extremos.

Para $n=0$ a identidade diz que $\int_{0}^{x}e^{is}ds=x+i\int_{0}^{x}(x-s)e^{is}ds$,
e o lado esquerdo é $(e^{ix}-1)/i$; multiplicando por $i$ e isolando $e^{ix}$,
\[
e^{ix}=1+ix+i^{2}\int_{0}^{x}(x-s)e^{is}ds ,
\]
que é o desenvolvimento de Taylor de ordem $1$ com resto integral: o
coeficiente de $x^{2}$ na série é $i^{2}/2!$, e a integração de $(x-s)$ produz
o $2$ que falta. Usando o caso $n=1$ e substituindo o valor da integral,
\[
e^{ix}=1+ix+\frac{i^{2}x^{2}}{2}+\frac{i^{3}}{2}\int_{0}^{x}(x-s)^{2}e^{is}ds .
\]
O padrão está claro. Se valer
\[
e^{ix}-\sum_{m=0}^{n}\frac{(ix)^{m}}{m!}=\frac{i^{n+1}}{n!}\int_{0}^{x}(x-s)^{n}e^{is}ds
\]
para um certo $n$, então a identidade aplicada ao lado direito fornece
\[
\frac{i^{n+1}}{n!}\int_{0}^{x}(x-s)^{n}e^{is}ds=\frac{i^{n+1}x^{n+1}}{(n+1)!}+\frac{i^{n+2}}{(n+1)!}\int_{0}^{x}(x-s)^{n+1}e^{is}ds ,
\]
já que $n!\,(n+1)=(n+1)!$. Isso é a mesma fórmula com $n$ trocado por $n+1$,
o que conclui a indução. Portanto, para todo $n\geqslant0$,
\[
e^{ix}-\sum_{m=0}^{n}\frac{(ix)^{m}}{m!}=\frac{i^{n+1}}{n!}\int_{0}^{x}(x-s)^{n}e^{is}ds .
\]
Toda a dificuldade está agora na integral do lado direito, que é o resto.

\medskip
\textbf{Segunda etapa: a cota boa para $|x|$ pequeno.} Como $|i^{n+1}|=1$ e
$|e^{is}|\leqslant1$, e mudando a variável para $u=|s|$ (o que também cobre
$x<0$),
\[
\Bigl|\frac{i^{n+1}}{n!}\int_{0}^{x}(x-s)^{n}e^{is}ds\Bigr|\leqslant\frac{1}{n!}\int_{0}^{|x|}(|x|-u)^{n}du=\frac{|x|^{n+1}}{(n+1)!} .
\]
O denominador vem daí: o $n!$ já estava na fórmula e a integral de $(x-s)^{n}$
em $(0,x)$ acrescenta o fator $n+1$. Esta é a cota do cálculo diferencial usual,
ótima para $|x|$ pequeno.

\medskip
\textbf{Terceira etapa: a cota boa para $|x|$ grande.}
Suponha $n\geqslant1$ e integre por partes outra vez, com $dv=(x-s)^{n-1}ds$:
\begin{align*}
\int_{0}^{x}(x-s)^{n-1}e^{is}ds
&=\frac{x^{n}}{n}+\frac{i}{n}\int_{0}^{x}(x-s)^{n}e^{is}ds ,
\\[0.3cm]
\frac{i}{n}\int_{0}^{x}(x-s)^{n}e^{is}ds
&=-\frac{x^{n}}{n}+\int_{0}^{x}(x-s)^{n-1}e^{is}ds .
\end{align*}
O ponto decisivo é que o termo de fronteira é ele próprio uma integral da
mesma família, $x^{n}/n=\int_{0}^{x}(x-s)^{n-1}ds$; substituindo, os dois
pedaços se fundem em uma única integral,
\[
\frac{i^{n+1}}{n!}\int_{0}^{x}(x-s)^{n}e^{is}ds=\frac{i^{n}}{(n-1)!}\int_{0}^{x}(x-s)^{n-1}\bigl(e^{is}-1\bigr)ds ,
\]
pois $\frac{i^{n}}{(n-1)!}\cdot\frac{i}{n}=\frac{i^{n+1}}{n!}$. É aqui que a
subtração de $1$ faz o seu trabalho: sem ela teríamos duas parcelas, uma do
termo de fronteira e outra da integral, e ao cotá-las separadamente voltaríamos
à ordem de grandeza $|x|^{n+1}$. Com o $-1$ dentro do parêntese resta uma só
integral, cujo integrando é uniformemente limitado, já que
$|e^{is}-1|\leqslant|e^{is}|+1\leqslant2$ para todo $s$. Essa limitação não
depende de $|x|$, e é isso que permite conservar o fator $|x|^{n}$:
\[
\Bigl|\frac{i^{n+1}}{n!}\int_{0}^{x}(x-s)^{n}e^{is}ds\Bigr|\leqslant\frac{2}{(n-1)!}\int_{0}^{|x|}(|x|-u)^{n-1}du=\frac{2|x|^{n}}{n!} .
\]

\medskip
\textbf{O caso $n=0$ e a conclusão.}
Para $n=0$ a segunda etapa já dá $|e^{ix}-1|\leqslant|x|$ e a terceira não se
aplica, pois divide por $n$; a segunda cota do mínimo é então $2$, que é a
estimativa trivial $|e^{ix}-1|\leqslant2$. Assim o enunciado vale para todo
$n\geqslant0$: para $n\geqslant1$ porque as duas últimas etapas fornecem ambas
as cotas, e vale a menor; para $n=0$ pela observação anterior.
\end{proof}

A transposição para funções características é imediata. Se
$\mathbb{E}|X|^{n}<\infty$, então
\[
\min\bigl(|tX|^{n+1},2|tX|^{n}\bigr)\leqslant2|t|^{n}|X|^{n},
\]
e o lado direito dessa desigualdade é integrável. Assim todos os termos abaixo
são finitos. Aplicando o lema a $x=tX$ e tomando esperanças, obtemos, para todo
$t\in\mathbb{R}$,
\begin{equation}\label{eq:cota-resto-taylor}
\begin{aligned}
\Bigl|\mathbb{E}\bigl[e^{itX}\bigr]
-\sum_{m=0}^{n}\mathbb{E}\Bigl[\frac{(itX)^{m}}{m!}\Bigr]\Bigr|
&\leqslant\mathbb{E}\Bigl[\Bigl|e^{itX}
-\sum_{m=0}^{n}\frac{(itX)^{m}}{m!}\Bigr|\Bigr]
\\[0.3cm]
&\leqslant\mathbb{E}\Bigl[\min\bigl(|tX|^{n+1},\,2|tX|^{n}\bigr)\Bigr].
\end{aligned}
\end{equation}
Primeiro o módulo entra na esperança, pela desigualdade triangular
(equivalentemente, pela desigualdade de Jensen aplicada a $|\cdot|$); depois se
usa a cota pontual do lema. Os denominadores foram descartados porque
$(n+1)!\geqslant1$ e $n!\geqslant1$, de modo que retirá-los só aumenta o lado
direito.

A cota \eqref{eq:cota-resto-taylor} vale para todo $t$ real, sem restrição de
tamanho, mas o interesse é o comportamento quando $t\to0$, e aí as duas
parcelas do mínimo competem: se $\mathbb{E}|X|^{n+1}<\infty$, a primeira dá
$O(|t|^{n+1})$; se dispomos apenas de $\mathbb{E}|X|^{n}<\infty$, a segunda dá
$O(|t|^{n})$. Em \Cref{thm:taylor-segunda-ordem} usaremos $n=2$: a primeira
alternativa exigiria o terceiro momento finito, e é a segunda que permite obter
a expansão de segunda ordem da função característica sob a hipótese de variância
finita apenas. Em resumo, o resto de Taylor de ordem $n$ é $O(|t|^{n+1})$ se
existirem $n+1$ momentos e $O(|t|^{n})$ se existirem apenas $n$.

\subsection{A expansão de segunda ordem e sua recíproca}\label{subsec:segunda-ordem}

O caso especial $n=2$ da cota \eqref{eq:cota-resto-taylor} será útil adiante, quando
tratarmos de somas de variáveis independentes. Com $n=2$, o polinômio de Taylor de
$\varphi_X$ na origem é
\[
1+it\,\mathbb{E}[X]+\frac{(it)^2}{2}\,\mathbb{E}[X^2]=1+it\,\mathbb{E}[X]-\frac{t^2}{2}\,\mathbb{E}[X^2],
\]
e o resto é majorado por $\mathbb{E}\bigl[\min\bigl(|tX|^3,2|tX|^2\bigr)\bigr]$. Como
$\mathbb{E}|X|^2<\infty$ implica $\mathbb{E}|X|<\infty$, os dois coeficientes são finitos.

\begin{theorem}[Expansão de segunda ordem]\label{thm:taylor-segunda-ordem}
Seja $X$ uma variável aleatória com função característica $\varphi_X$.
Se $\mathbb{E}|X|^2<\infty$, então
\[
\varphi_X(t)=1+it\,\mathbb{E}[X]-\frac{t^2}{2}\,\mathbb{E}[X^2]+o(t^2),\qquad t\to0.
\]
\end{theorem}

\begin{proof}
Seja $R_2(t)\equiv\varphi_X(t)-1-it\,\mathbb{E}[X]+\tfrac{t^2}{2}\mathbb{E}[X^2]$ o resto, de
modo que, por \eqref{eq:cota-resto-taylor},
$|R_2(t)|\leqslant\mathbb{E}\bigl[\min\bigl(|tX|^3,2|tX|^2\bigr)\bigr]$. Vamos isolar o fator
$t^2$ dentro da esperança, para que a divisão por $t^2$ apareça explicitamente. Como
$|tx|^3=t^2|t||x|^3$ e $2|tx|^2=2t^2|x|^2$, dividindo ambos os lados por $t^2>0$ obtemos
\begin{equation}\label{eq:resto-dividido}
\frac{|R_2(t)|}{t^2}\leqslant\mathbb{E}\bigl[\min\bigl(|t|\,|X|^3,2|X|^2\bigr)\bigr].
\end{equation}
Basta mostrar que o lado direito tende a zero, o que é a afirmação $R_2(t)=o(t^2)$. Ponha
$Y_t\equiv\min\bigl(|t||X|^3,2|X|^2\bigr)$. Temos $0\leqslant Y_t\leqslant2|X|^2$, e
$2|X|^2$ é integrável por hipótese, logo $(Y_t)$ tem dominante integrável que não depende de
$t$. Para a convergência pontual, fixe $\omega$ com $|X(\omega)|<\infty$, o que vale quase
certamente porque $\mathbb{E}|X|^2$ é finito; quando $t\to0$, o primeiro argumento do mínimo
tende a $0$ e o segundo é a constante $2|X(\omega)|^2$, de sorte que $Y_t(\omega)\to0$. Pelo
teorema da convergência dominada, $\mathbb{E}[Y_t]\to0$; como $|R_2(t)|/t^2\geqslant0$, segue
de \eqref{eq:resto-dividido} que
\[
0\leqslant\limsup_{t\to0}\frac{|R_2(t)|}{t^2}\leqslant\lim_{t\to0}\mathbb{E}[Y_t]=0,
\]
isto é, $|R_2(t)|/t^2\to0$ e portanto $R_2(t)=o(t^2)$.
\end{proof}

\begin{remark}[Por que a cota com o mínimo]\label{rem:cota-minimo}
O ponto da cota \eqref{eq:cota-resto-taylor}, que envolve o mínimo de dois termos em vez do
primeiro termo que resultaria de uma aplicação ingênua da série de Taylor, é que assim a
conclusão de \Cref{thm:taylor-segunda-ordem} vale sob a hipótese $\mathbb{E}|X|^2<\infty$,
sem exigir $\mathbb{E}|X|^3<\infty$.

A cota ingênua para $n=2$, $\bigl|e^{ix}-1-ix+x^2/2\bigr|\leqslant|x|^3/3!$, majora o resto
por $\tfrac{1}{6}|t|^3\,\mathbb{E}|X|^3$ e, após a divisão por $t^2$, deixa
$\tfrac{1}{6}|t|\,\mathbb{E}|X|^3$, que só tende a zero se o terceiro momento absoluto for
finito.

O mínimo contorna essa exigência porque trata de modo diferente as duas regiões
determinadas por $|tX|$. Se $|tX|\leqslant2$, então $|tX|^3\leqslant2|tX|^2$ e o mínimo é o
termo cúbico, o termo fino, útil justamente onde $|tX|$ é pequeno; como ele carrega o fator
adicional $|t|$, decai quando $t\to0$, mesmo que $\mathbb{E}|X|^3$ seja infinito. Se
$|tX|>2$, o mínimo é $2|tX|^2$, que não envolve $|X|^3$: nessa região, onde valores grandes
de $|X|$ custariam integrabilidade, a cota é proporcional a $t^2|X|^2$, e o fator $t^2$ já
está fora da esperança, pronto para ser cancelado na divisão. Em suma, o fator extra $|t|$
da parte cúbica absorve o que seria uma potência maior de $|X|$, e a parte limitada fica
controlada por $|X|^2$.
\end{remark}

O próximo resultado mostra que a existência da segunda derivada de $\varphi_X$ na origem
implica a existência do segundo momento, isto é, ele é a recíproca de
\Cref{thm:taylor-segunda-ordem}. Como $\varphi_X$ assume valores complexos, a derivada deve
ser entendida no sentido da diferença segunda simétrica, expressão que aparece na hipótese.

\begin{theorem}[Segundo momento e diferença segunda]\label{thm:segundo-momento-finito}
Seja $X$ uma variável aleatória com função característica $\varphi_X$. Se
\[
\limsup_{h\downarrow0}\frac{\varphi_X(h)-2\varphi_X(0)+\varphi_X(-h)}{h^2}>-\infty,
\]
então $\mathbb{E}|X|^2<\infty$.
\end{theorem}

\begin{proof}
Seja $F$ a função de distribuição de $X$, de modo que $\varphi_X(t)=\int e^{itx}\,dF(x)$. Pela
linearidade da integral, aplicada ao integrando limitado,
\[
\frac{\varphi_X(h)-2\varphi_X(0)+\varphi_X(-h)}{h^2}=\int_{\mathbb{R}}\frac{e^{ihx}-2+e^{-ihx}}{h^2}\,dF(x).
\]
Como $e^{ihx}+e^{-ihx}=2\cos(hx)$, o integrando vale
\begin{equation}\label{eq:integrando-segunda-diferenca}
\frac{e^{ihx}-2+e^{-ihx}}{h^2}=-\frac{2\bigl(1-\cos(hx)\bigr)}{h^2}\leqslant0,
\end{equation}
pois $1-\cos(hx)\geqslant0$; e o limite fundamental $\lim_{u\to0}(1-\cos u)/u^2=1/2$, com
$u=hx$, mostra que $2\bigl(1-\cos(hx)\bigr)/h^2\to x^2$ quando $h\to0$, para cada $x$ fixado.

O sinal em \eqref{eq:integrando-segunda-diferenca} é o que autoriza usar o lema de Fatou na
forma que fornece cota inferior, pois Fatou se aplica a funções não negativas: se
$g_h\geqslant0$, então $\int\liminf_{h\to0}g_h\,dF\leqslant\liminf_{h\to0}\int g_h\,dF$.
Aplicado diretamente ao integrando negativo, o lema daria a desigualdade no sentido oposto,
inútil aqui. Tomando então $g_h(x)\equiv2\bigl(1-\cos(hx)\bigr)/h^2\geqslant0$, o simétrico
do integrando original, e lembrando que $2\varphi_X(0)=2$ e que
$\varphi_X(h)+\varphi_X(-h)=2\int\cos(hx)\,dF(x)$, obtemos
\begin{align*}
\int_{\mathbb{R}}x^2\,dF(x)
&= \int_{\mathbb{R}}\lim_{h\to0}\frac{2\bigl(1-\cos(hx)\bigr)}{h^2}\,dF(x)
\\[0.3cm]
&\leqslant \liminf_{h\to0}\int_{\mathbb{R}}\frac{2\bigl(1-\cos(hx)\bigr)}{h^2}\,dF(x)
\\[0.3cm]
&= \liminf_{h\to0}\frac{2\varphi_X(0)-\varphi_X(h)-\varphi_X(-h)}{h^2}
\\[0.3cm]
&= -\limsup_{h\to0}\frac{\varphi_X(h)-2\varphi_X(0)+\varphi_X(-h)}{h^2}<\infty,
\end{align*}
sendo a última desigualdade exatamente a hipótese. Como $x\mapsto x^2$ é não
negativa, a integral $\int x^2\,dF(x)$ é o segundo momento $\mathbb{E}[X^2]$, e
a cadeia mostra que esse número é finito, ou seja, que
$\mathbb{E}|X|^2<\infty$.
\end{proof}

Combinando \Cref{thm:taylor-segunda-ordem} e \Cref{thm:segundo-momento-finito}, a finitude
do segundo momento é equivalente à existência da segunda derivada de $\varphi_X$ na origem, no
sentido da diferença segunda simétrica. Se $\mathbb{E}|X|^2<\infty$, a expansão de segunda
ordem diz que $\bigl(\varphi_X(h)-2\varphi_X(0)+\varphi_X(-h)\bigr)/h^2$ converge para
$-\mathbb{E}[X^2]$; reciprocamente, a hipótese de \Cref{thm:segundo-momento-finito} garante
o segundo momento finito. Note que não se pede a existência do limite, apenas que o limite
superior não seja $-\infty$.

\section{Exercícios Propostos}\label{sec:exercicios}

Os exercícios a seguir acompanham as seções anteriores na mesma ordem em que
aparecem no texto.

\begin{exercise}[Parte real e módulo ao quadrado de uma função característica]%
\label{exercise:parte-real-modulo}
Seja $X$ uma variável aleatória com função característica $\varphi_X$. Mostre que
$\operatorname{Re}\varphi_X$ e $|\varphi_X|^2$ também são funções características.

\emph{Sugestão:} em cada caso, exiba uma variável aleatória cuja função
característica seja a função pedida, trabalhando com cópias independentes de $X$
e, no primeiro caso, com um sinal de Rademacher independente de $X$. Use a
independência para fatorar as esperanças.
\end{exercise}

\begin{exercise}[Inversão para massas pontuais]\label{exercise:massa-pontual}
Seja $X$ uma variável aleatória com função característica $\varphi_X$.
\begin{enumerate}
\item[(i)] Imitando a prova do \Cref{thm:inversao}, mostre que, para todo
$a\in\mathbb{R}$,
\[
\mathbb{P}(X=a)
=\lim_{T\to\infty}\frac{1}{2T}\int_{-T}^{T}e^{-ita}\,\varphi_X(t)\,dt .
\]

\item[(ii)] Suponha que $\mathbb{P}(X\in h\mathbb{Z})=1$ para algum $h>0$.
Mostre que $\varphi_X$ é periódica de período $2\pi/h$, isto é,
\[
\varphi_X\Bigl(t+\frac{2\pi}{h}\Bigr) = \varphi_X(t), \qquad t\in\mathbb{R},
\]
e conclua a fórmula de inversão discreta
\[
\mathbb{P}(X=x) = \frac{h}{2\pi}\int_{-\pi/h}^{\pi/h}e^{-itx}\,\varphi_X(t)\,dt,
\qquad x\in h\mathbb{Z}.
\]

\item[(iii)] Se $X=Y+b$ com $b\in\mathbb{R}$, então
$\mathbb{E}[e^{itX}] = e^{itb}\,\mathbb{E}[e^{itY}]$. Deduza daí que, se
$\mathbb{P}(X\in b+h\mathbb{Z})=1$, a fórmula de inversão de (ii) continua
válida, na mesma forma, para $x\in b+h\mathbb{Z}$.
\end{enumerate}

\emph{Sugestão}: em (i), troque esperança e integral como na prova do
\Cref{thm:inversao} e examine o limite do núcleo que aparece quando
$T\to\infty$. Em (ii), use que $e^{2\pi iX/h}=1$ quase certamente para obter a
periodicidade e, depois, escreva o limite de (i) como uma média de Cesàro de uma
função periódica de período $2\pi/h$, que converge para a média sobre um
período. Em (iii), aplique (ii) a $Y=X-b$ e reúna os fatores exponenciais.
\end{exercise}

\begin{exercise}[Média quadrática da função característica]%
\label{exercise:media-quadrado}
Sejam $X$ e $Y$ independentes e com a mesma distribuição, de modo que
$\varphi_X=\varphi_Y$.
Aplicando o \Cref{exercise:massa-pontual} à variável aleatória $X-Y$, conclua que
\[
\lim_{T\to\infty}\frac{1}{2T}\int_{-T}^{T}|\varphi_X(t)|^2\,dt
= \mathbb{P}(X-Y=0)
= \sum_{x}\mathbb{P}(X=x)^2 .
\]
A soma se estende aos valores $x$ com $\mathbb{P}(X=x)>0$, que são no máximo
enumeráveis. No último passo, use a independência de $X$ e $Y$ para expressar
$\mathbb{P}(X=Y)$ em termos dessas probabilidades pontuais.
\end{exercise}

\begin{remark}\label{rem:sem-massa-pontual}
Seja $X$ uma variável aleatória com função característica $\varphi_X$.
A conclusão do \Cref{exercise:media-quadrado} implica que, se
$\varphi_X(t)\to0$ quando $t\to\infty$, então $\mathbb{P}(X=x)=0$ para todo
$x\in\mathbb{R}$. De fato, nesse caso $|\varphi_X(t)|^2\to0$ quando
$|t|\to\infty$, pois $\varphi_X(-t)=\overline{\varphi_X(t)}$. Como o integrando é
limitado por $1$, sua média em $[-T,T]$ também tende a zero. Logo
$\sum_x\mathbb{P}(X=x)^2=0$ e cada parcela se anula.

O \Cref{exercise:cantor} mostra que a recíproca é falsa: uma distribuição pode
não ter massas pontuais e, ainda assim, sua função característica não tender a
zero no infinito. Por outro lado, um resultado conhecido como lema de
Riemann--Lebesgue garante que, se $X$ tem densidade $f$, então
$\varphi_X(t)=\int_{\mathbb{R}}f(x)e^{itx}\,dx$ tende a zero quando
$|t|\to\infty$.
\end{remark}

\begin{exercise}[Densidade com função característica não integrável]%
\label{exercise:densidade-integral-divergente}
Dê um exemplo de variável aleatória $X$ que tenha densidade e cuja função
característica $\varphi_X$ satisfaça
\[
\int_{-\infty}^{\infty}|\varphi_X(t)|\,dt = \infty .
\]
\emph{Sugestão}: dois dos exemplos anteriores, a saber, o do
\Cref{ex:uniforme} e o do \Cref{ex:exponencial}, têm essa propriedade. Para a
exponencial, compare $|\varphi_X(t)|$ com $1/|t|$ quando $|t|$ é grande. Para a
uniforme em $(-1,1)$, examine os intervalos em que $|\sin t|$ está afastado de
zero.
\end{exercise}

\begin{exercise}[Soma de variáveis uniformes independentes]%
\label{exercise:soma-uniformes}
Suponha que $X_1,\ldots,X_n$ sejam independentes e uniformemente distribuídas no
intervalo $(-1,1)$. Mostre que, para $n\geqslant 2$, a soma
$X_1+\cdots+X_n$ tem densidade
\[
f(x) = \frac{1}{\pi}\int_{0}^{\infty}\Bigl(\frac{\sin t}{t}\Bigr)^{n}
\cos(tx)\,dt .
\]
Embora não seja evidente a partir da fórmula, $f$ é um polinômio em cada
intervalo $(k,k+1)$ com $k\in\mathbb{Z}$ e se anula fora de $[-n,n]$.

\emph{Sugestão}: obtenha a função característica da soma com o
\Cref{thm:fc-soma-independentes}, aplique a fórmula de inversão na forma do
\Cref{thm:densidade-fc-integravel}, explorando a paridade do integrando, e, para
a afirmação final, lembre que a densidade da soma é a convolução de $n$
densidades uniformes, logo é polinomial por partes em cada intervalo unitário.
\end{exercise}

\begin{exercise}[Médias de variáveis com distribuição de Cauchy]%
\label{exercise:cauchy-media}
Use o resultado do \Cref{ex:cauchy} para concluir que, se $X_1,X_2,\ldots$ são
independentes e têm distribuição de Cauchy, então, para todo $n\geqslant 1$, a
média $(X_1+\cdots+X_n)/n$ tem a mesma distribuição que $X_1$.

\emph{Sugestão}: calcule a função característica da média usando a
independência e compare-a com a função característica de $X_1$; duas variáveis
aleatórias com a mesma função característica têm a mesma distribuição, pela
fórmula de inversão (\Cref{thm:inversao}).
\end{exercise}

\begin{exercise}[Variâncias de normais convergentes]
\label{exercise:normal-variancia}
Suponha que $X_n\Rightarrow X$ e que, para cada $n\geqslant1$, $X_n$ tenha
distribuição normal de média $0$ e variância $\sigma_n^2$, com
$0\leqslant\sigma_n^2<\infty$. Mostre que existe $\sigma^2\in[0,\infty)$ tal que
$\sigma_n^2\to\sigma^2$. O limite pode ser $0$: nesse caso $X\equiv0$, a
distribuição degenerada em $0$, que interpretamos como a normal de variância
$0$. Nada se supõe de antemão sobre a convergência da sequência numérica
$(\sigma_n^2)$, nem sobre a natureza da distribuição limite.

\emph{Sugestão:} por \Cref{thm:continuidade},
$e^{-\sigma_n^2t^2/2}=\varphi_{X_n}(t)\to\varphi_X(t)$ para cada
$t\in\mathbb{R}$. Tome uma subsequência com
$\sigma_{n_k}^2\to\sigma^2\in[0,\infty]$, com a convenção $e^{-\infty}=0$; ao
longo dela, $\varphi_X(t)=e^{-\sigma^2t^2/2}$ para todo $t$. Se
$\sigma^2=\infty$, teríamos $\varphi_X(t)=0$ para todo $t\neq0$, contra
$\varphi_X(0)=1$ e a continuidade de $\varphi_X$ na origem. Logo
$\sigma^2<\infty$ e o limite não depende da subsequência escolhida; pelo
critério de subsequências para convergência de números reais,
$\sigma_n^2\to\sigma^2$.
\end{exercise}

\begin{exercise}[Somas de pares independentes convergentes]
\label{exercise:soma-independentes-convergencia}
Para cada $n\in\{1,2,\dots\}\cup\{\infty\}$, suponha que $X_n$ e $Y_n$ sejam
independentes, isto é, que o par $(X_n,Y_n)$ seja formado por variáveis
independentes; o mesmo vale para o par limite. Não se supõe independência entre
variáveis de índices diferentes. Suponha ainda que $X_n\Rightarrow X_\infty$ e
$Y_n\Rightarrow Y_\infty$. Mostre que
$X_n+Y_n\Rightarrow X_\infty+Y_\infty$.

\emph{Sugestão:} pela independência dentro de cada par e por
\Cref{thm:fc-soma-independentes},
\[
\varphi_{X_n+Y_n}(t)=\varphi_{X_n}(t)\,\varphi_{Y_n}(t),
\qquad t\in\mathbb{R},
\]
para todo $n$, inclusive $n=\infty$. Por \Cref{thm:continuidade} aplicado às
duas sequências, o produto converge para $\varphi_{X_\infty}(t)\varphi_{Y_\infty}(t)$,
que é a função característica de $X_\infty+Y_\infty$ e é contínua em $0$. Use
\Cref{thm:continuidade} na direção contrária.
\end{exercise}

\begin{exercise}[Função característica de um produto infinito]
\label{exercise:produto-infinito-qc}
Sejam $X_1,X_2,\dots$ variáveis aleatórias independentes, $S_n=X_1+\cdots+X_n$
e $\varphi_{X_j}$ a função característica de $X_j$. Suponha que $S_n\to S_\infty$
quase certamente. Mostre que
\[
  \varphi_{S_\infty}(t)=\prod_{j=1}^{\infty}\varphi_{X_j}(t),
  \qquad t\in\mathbb{R},
\]
onde o produto infinito é, por definição, o limite dos produtos parciais
$\prod_{j=1}^{n}\varphi_{X_j}(t)$, cuja existência é parte da conclusão. Nenhuma
hipótese de integrabilidade é feita sobre as $X_j$.

\emph{Sugestão:} a convergência quase certa implica convergência em
distribuição, como vimos nas notas de aula sobre modos de convergência; logo
$S_n\Rightarrow S_\infty$ e, por \Cref{thm:continuidade},
$\varphi_{S_n}(t)\to\varphi_{S_\infty}(t)$ para todo $t$. Pela independência,
uma indução sobre $n$ a partir de \Cref{thm:fc-soma-independentes} dá
$\varphi_{S_n}(t)=\varphi_{X_1}(t)\cdots\varphi_{X_n}(t)$, e basta identificar o limite.
\end{exercise}

\begin{exercise}[O produto de cossenos]
\label{exercise:produto-cossenos}
Iterando a identidade $\sin t=2\sin(t/2)\cos(t/2)$ obtemos
\[
  \frac{\sin t}{t}=\prod_{m=1}^{\infty}\cos\!\left(\frac{t}{2^{m}}\right),
  \qquad t\neq0.
\]
Prove essa identidade interpretando cada um dos lados como função característica
de uma mesma distribuição. Para isso: (i) mostre por indução que
\[
\sin t=2^{n}\sin\bigl(t/2^{n}\bigr)\prod_{m=1}^{n}\cos\bigl(t/2^{m}\bigr)
\]
e conclua que $(\sin t)/t$ é o limite dos produtos parciais, já que
$2^{n}\sin(t/2^{n})\to t$ quando $n\to\infty$; (ii) verifique que
$t\mapsto(\sin t)/t$, com o valor $1$ em $0$, é a função característica da
distribuição uniforme em $[-1,1]$; (iii) verifique que $t\mapsto\cos(t/2^{m})$
é a função característica de $Y_m$, que assume os valores $\pm2^{-m}$ com
probabilidade $1/2$ cada; (iv) aplique \Cref{exercise:produto-infinito-qc} às
variáveis independentes $Y_1,Y_2,\dots$ e mostre que $\sum_{m\geqslant1}Y_m$ é
uniforme em $[-1,1]$. No último passo, note que a série converge absolutamente,
pois é majorada por $\sum_{m\geqslant1}2^{-m}=1$, e que, escrevendo
$Y_m=2^{-m}(2U_m-1)$ com $U_m$ de Bernoulli simétrica, a soma coincide em
distribuição com $2U-1$, onde $U=\sum_{m\geqslant1}U_m2^{-m}$ é uniforme em
$[0,1]$, por ser a série dos dígitos da expansão binária. A unicidade da função
característica, dada por \Cref{thm:inversao}, fecha o argumento.
\end{exercise}

\begin{exercise}[A distribuição de Cantor]
\label{exercise:cantor}
Sejam $X_1,X_2,\dots$ independentes, cada uma assumindo os valores $0$ e $1$
com probabilidade $1/2$. Defina
\[
  X=2\sum_{j\geqslant1}\frac{X_j}{3^{j}}.
\]
A série converge quase certamente, e mesmo absolutamente, pois a soma dos
módulos de seus termos é limitada por $1$. Logo $X$ está bem definida, toma
valores em $[0,1]$ e tem a distribuição de Cantor, suportada no conjunto
ternário obtido removendo os terços médios. Calcule a função característica
$\varphi_X$ de $X$ e conclua que
\[
  \varphi_X(t)=e^{it/2}\prod_{j=1}^{\infty}\cos\!\left(\frac{t}{3^{j}}\right),
  \qquad t\in\mathbb{R}.
\]
Mostre em seguida, a partir da autossemelhança
$X\overset{d}{=}(2X_1+X')/3$, em que $X'$ é independente de $X_1$ e tem a mesma
distribuição de $X$, que $\varphi_X(3t)=e^{it}\cos t\,\varphi_X(t)$, e conclua que
$\varphi_X$ assume o mesmo valor nos pontos $t=3^{k}\pi$, isto é,
$\varphi_X(3^{k}\pi)=\varphi_X(\pi)$ para todo $k=0,1,2,\dots$. Interprete esse fato
à luz do decaimento das funções características de distribuições com densidade.

\emph{Sugestão:} cada parcela $2X_j/3^{j}$ assume $0$ e $2/3^{j}$ com
probabilidade $1/2$ e tem função característica
$\tfrac12\bigl(1+e^{2it/3^{j}}\bigr)=e^{it/3^{j}}\cos(t/3^{j})$; use
\Cref{exercise:produto-infinito-qc} para o produto infinito e
$\sum_{j\geqslant1}3^{-j}=1/2$ para recolher as exponenciais. Na recorrência,
note que $e^{i3^{k}\pi}\cos(3^{k}\pi)=1$ para todo $k\geqslant1$.
\end{exercise}

\begin{exercise}[Momentos da normal padrão]
\label{exercise:momentos-normal}
Seja $X$ uma variável aleatória com distribuição normal padrão. Use
\Cref{thm:derivadas-momentos} e o desenvolvimento em série de potências de
$e^{-t^{2}/2}$ para mostrar que, para todo $n\geqslant1$,
\[
  \mathbb{E}\bigl[X^{2n}\bigr]=\frac{(2n)!}{2^{n}\,n!}
  =(2n-1)(2n-3)\cdots3\cdot1\equiv(2n-1)!!,
\]
onde $(2n-1)!!$ é, por definição, o produto dos inteiros positivos ímpares que
não excedem $2n-1$, com a convenção $(-1)!!=1$. Em particular,
$\mathbb{E}[X^{2}]=1$ e $\mathbb{E}[X^{4}]=3$. Mostre também que os momentos de
ordem ímpar se anulam, por simetria da distribuição:
$\mathbb{E}[X^{2n-1}]=0$ para $n\geqslant1$.

\emph{Sugestão:} a função característica da normal padrão é
$\varphi_X(t)=e^{-t^{2}/2}=\sum_{n\geqslant0}\frac{(-1)^{n}}{2^{n}n!}t^{2n}$, e
\Cref{thm:derivadas-momentos} dá
$\varphi_X^{(2n)}(0)=i^{2n}\mathbb{E}[X^{2n}]=(-1)^{n}\mathbb{E}[X^{2n}]$;
compare os coeficientes das duas expressões de $\varphi_X$. Para os momentos
ímpares, use $X\overset{d}{=}-X$.
\end{exercise}

\begin{exercise}[Equicontinuidade das funções características]\label{exercise:equicontinuidade}
Dizemos que uma família de variáveis aleatórias $\{X_i:i\in I\}$ é
\emph{condensada} se, para todo $\varepsilon>0$, existe $M>0$ tal que
\[
\mathbb{P}(|X_i|>M)<\varepsilon
\qquad\text{para todo }i\in I.
\]
Denotando por $\varphi_{X_i}$ a função característica de $X_i$, dizemos que a
família $\{\varphi_{X_i}:i\in I\}$ é \emph{equicontínua} se, para todo
$\varepsilon>0$, existe $\delta>0$ tal que $|h|<\delta$ implica
$|\varphi_{X_i}(t+h)-\varphi_{X_i}(t)|<\varepsilon$ para todos $i\in I$ e
$t\in\mathbb{R}$.

\begin{enumerate}
\item[(i)] Suponha que $\{X_i:i\in I\}$ seja condensada e mostre que
$\{\varphi_{X_i}:i\in I\}$ é equicontínua.

\emph{Sugestão:} pelo item~(d) do \Cref{thm:propriedades-fc},
\[
|\varphi_{X_i}(t+h)-\varphi_{X_i}(t)|
\leqslant\mathbb{E}\bigl[|e^{ihX_i}-1|\bigr].
\]
Use $|e^{i\theta}-1|\leqslant|\theta|\wedge2$, obtida em
\eqref{eq:cota-resto-taylor}, para mostrar que, para todo $M>0$,
\[
|\varphi_{X_i}(t+h)-\varphi_{X_i}(t)|
\leqslant |h|M+2\mathbb{P}(|X_i|>M).
\]
Escolha primeiro $M$ e depois $\delta$ em função de $\varepsilon$.

\item[(ii)] Suponha que $X_n\Rightarrow X_\infty$ e denote por
$\varphi_{X_\infty}$ a função característica de $X_\infty$. Mostre que a
família $\{X_n:n\in\mathbb{N}\}\cup\{X_\infty\}$ é condensada e conclua que
$\varphi_{X_n}\to\varphi_{X_\infty}$ uniformemente sobre cada intervalo limitado.

\emph{Sugestão:} para obter a condensação, escolha $M$ de modo que $-M$ e $M$
sejam pontos de continuidade de $F_{X_\infty}$ e que
$\mathbb{P}(|X_\infty|\geqslant M)$ seja pequena. Use a convergência das
funções de distribuição e, depois, aumente $M$ para incluir os finitos
índices iniciais. Aplique o item~(i) e o \Cref{thm:continuidade}. Dado um
intervalo limitado, cubra-o por um número finito de intervalos de raio
$\delta$, centrados em $t_1,\ldots,t_k$, e escreva
\[
\begin{aligned}
|\varphi_{X_n}(t)-\varphi_{X_\infty}(t)|
&\leqslant|\varphi_{X_n}(t)-\varphi_{X_n}(t_j)|
+|\varphi_{X_n}(t_j)-\varphi_{X_\infty}(t_j)|\\
&\quad+|\varphi_{X_\infty}(t_j)-\varphi_{X_\infty}(t)|.
\end{aligned}
\]
Use a equicontinuidade na primeira e na terceira parcelas e a convergência
pontual em $t_j$ na parcela do meio.

\item[(iii)] Dê um exemplo em que $X_n\Rightarrow X_\infty$, mas
$\varphi_{X_n}$ não converge para $\varphi_{X_\infty}$ uniformemente em
$\mathbb{R}$.

\emph{Sugestão:} tome $X_n\equiv1/n$ e $X_\infty\equiv0$ e observe que
$\sup_{t\in\mathbb{R}}|e^{it/n}-1|=2$.
\end{enumerate}
\end{exercise}

\begin{exercise}[Lei fraca e existência de $\varphi_{X_1}'(0)$]\label{exercise:pitman}
Sejam $X_1,X_2,\dots$ independentes e identicamente distribuídas, com função
característica $\varphi_{X_1}$, e $S_n=X_1+\cdots+X_n$.
\begin{enumerate}
\item[(i)] Se $\varphi_{X_1}'(0)=ia$, com $a\in\mathbb{R}$, mostre que
$S_n/n\to a$ em probabilidade. Como $\varphi_{S_n/n}(t)=\varphi_{X_1}(t/n)^n$ e, para
$t$ fixo,
\[
n\bigl(\varphi_{X_1}(t/n)-1\bigr)
=t\cdot\frac{\varphi_{X_1}(t/n)-1}{t/n}\longrightarrow iat,
\]
donde, com $\epsilon_n\to0$,
\[
\varphi_{X_1}(t/n)^n
=\Bigl(1+\frac{iat+\epsilon_n}{n}\Bigr)^{\!n}\longrightarrow e^{iat}.
\]
O \Cref{thm:continuidade} fornece então $S_n/n\Rightarrow a$, e a convergência
em distribuição para a constante $a$ equivale à convergência em probabilidade.

\item[(ii)] Reciprocamente, se $S_n/n\to a$ em probabilidade, mostre que
$\varphi_{X_1}(t/n)^n\to e^{iat}$ quando $n\to\infty$ pelos inteiros, com $t$ fixo.
Basta notar que $e^{itS_n/n}\to e^{ita}$ em probabilidade e que
$\bigl|e^{itS_n/n}\bigr|\leqslant1$, de modo que a convergência também se dá em
$L^1$.

\item[(iii)] Use o item~(ii) e a convergência uniforme sobre compactos do
\Cref{exercise:equicontinuidade}(ii) para mostrar que $(\varphi_{X_1}(h)-1)/h\to ia$
quando $h\to0$ pelos reais positivos.
\emph{Sugestão:} pela independência, a função característica de $S_n/n$ é
$\varphi_{S_n/n}(s)=\varphi_{X_1}(s/n)^n$. Pelo item~(ii),
$\varphi_{S_n/n}(s)\to e^{ias}$ para cada $s$, e o limite é
contínuo; logo o \Cref{exercise:equicontinuidade}(ii) dá convergência uniforme
sobre compactos. Fixe $\tau>0$ pequeno, com $|a|\tau<\pi/2$ e $\varphi_{X_1}(u)\neq0$
para $|u|\leqslant\tau$, de modo que o logaritmo principal esteja definido. Da
convergência uniforme em $[-\tau,\tau]$ segue que
$n\log\varphi_{X_1}(s/n)=\log\varphi_{S_n/n}(s)\to ias$ uniformemente nesse intervalo; como
$\log(1+z)=z+O(z^2)$ quando $z\to0$, obtemos
$n\bigl(\varphi_{X_1}(s/n)-1\bigr)\to ias$ uniformemente em $|s|\leqslant\tau$. Tome
agora $h\downarrow0$ e ponha $n=\lfloor\tau/h\rfloor$, de modo que
$s_n\equiv nh\to\tau$ e $h=s_n/n$. Então $n\bigl(\varphi_{X_1}(h)-1\bigr)\to ia\tau$ e,
dividindo por $nh\to\tau$, concluímos que $(\varphi_{X_1}(h)-1)/h\to ia$.
\end{enumerate}
Conclua que a lei fraca dos grandes números vale se, e somente se, $\varphi_{X_1}'(0)$
existe, e que nesse caso $\varphi_{X_1}'(0)=ia$.\footnote{%
O resultado deve-se a E.~J.~G.~Pitman (1956); a necessidade de demonstrar o
item~(iii) foi observada por John Walsh.}
\end{exercise}

\begin{remark}
Combinando o exercício anterior com o exemplo clássico de uma variável
aleatória $X$ com média infinita cuja função característica tem derivada na
origem, vê-se que a existência de $\varphi_X'(0)$ não garante
$\mathbb{E}|X|<\infty$. Um exemplo é uma variável aleatória simétrica com
caudas do tipo $c/|x|^2$, isto é, $\mathbb{P}(|X|>x)\sim c\,x^{-2}$ para $x$
grande: sua média diverge, pois
$\mathbb{E}|X|=\int_0^\infty\mathbb{P}(|X|>x)\,dx=\infty$, e no entanto
$\varphi_X(t)-1=O\bigl(t^2\log(1/|t|)\bigr)$ quando $t\to0$, de modo que
$(\varphi_X(t)-1)/t\to0$ e $\varphi_X'(0)=0$ existe.
\end{remark}

\begin{exercise}[Identidade para o primeiro momento]\label{exercise:primeiro-momento}
Seja $X$ uma variável aleatória com função de distribuição $F_X$ e função
característica $\varphi_X$. Mostre que
\[
  2\int_0^\infty\frac{1-\operatorname{Re}\varphi_X(t)}{\pi t^2}\,dt=\int|y|\,dF_X(y),
\]
com a igualdade entendida em $[0,+\infty]$. Sugestão: escreva
\[
1-\operatorname{Re}\varphi_X(t)
=\int\bigl(1-\cos(ty)\bigr)\,dF_X(y)
\]
e recorde, do \Cref{ex:distribuicao-polya}, que a densidade de Pólya
$\frac{1-\cos x}{\pi x^2}$ integra $1$ em $\mathbb{R}$; sendo ela par,
\[
  \int_0^\infty\frac{1-\cos x}{\pi x^2}\,dx=\frac12 .
\]
Com a mudança de variáveis $x=|y|t$ em
$\int_0^\infty\frac{1-\cos(ty)}{\pi t^2}\,dt$ obtém-se $\frac{|y|}{2}$, e o
teorema de Tonelli, já que o integrando é não negativo, conclui a troca das
integrais.
\end{exercise}

\begin{exercise}[Segunda derivada nula na origem]\label{exercise:segunda-derivada-nula}
Seja $X$ uma variável aleatória com função de distribuição $F_X$ e função
característica $\varphi_X$. Suponha que
\[
  \lim_{t\downarrow0}\frac{\varphi_X(t)-1}{t^2}=c>-\infty .
\]
Mostre que $\mathbb{E}[X]=0$ e que $\mathbb{E}|X|^2=-2c<\infty$. Em
particular, se $\varphi_X(t)=1+o(t^2)$ quando $t\to0$, então $\varphi_X(t)\equiv1$
e $X=0$ quase certamente.

\emph{Sugestão:} da hipótese segue que
$\bigl(1-\operatorname{Re}\varphi_X(t)\bigr)/t^2\leqslant C$ para
$t\in(0,\delta]$. Ora,
\[
\frac{1-\operatorname{Re}\varphi_X(t)}{t^2}
=\int_{\mathbb{R}}\frac{1-\cos(tx)}{t^2}\,dF_X(x),
\]
e o integrando é não negativo e converge para $x^2/2$ quando $t\downarrow0$,
para cada $x$ fixado. Aplicando o lema de Fatou ao longo de uma sequência
$t_k\downarrow0$ e depois o teorema da convergência monótona, obtemos
$\mathbb{E}[X^2]\leqslant2C<\infty$. Com o segundo momento finito, a expansão
de segunda ordem de \Cref{thm:taylor-segunda-ordem} dá
$\varphi_X(t)=1+it\,\mathbb{E}[X]-\frac{t^2}{2}\mathbb{E}[X^2]+o(t^2)$; divida
por $t^2$ e faça $t\downarrow0$ para concluir que $\mathbb{E}[X]=0$ e que
$c=-\frac{1}{2}\mathbb{E}[X^2]$.
\end{exercise}

\begin{exercise}[Convergência para zero]\label{exercise:convergencia-para-zero}
Sejam $Y_1,Y_2,\ldots$ variáveis aleatórias com funções características
$\varphi_{Y_n}$. Mostre que $Y_n\Rightarrow0$ se, e somente se, existe $\delta>0$
tal que $\varphi_{Y_n}(t)\to1$ para todo $t$ com $|t|\leqslant\delta$.

\emph{Sugestão:} numa das direções, use o \Cref{thm:continuidade}. Na outra,
suponha a convergência apenas em $[-\delta,\delta]$. Pela cota
\eqref{eq:cota-concentracao}, com $u=\delta$,
\[
\mathbb{P}\bigl(|Y_n|\geqslant2/\delta\bigr)
\leqslant\frac{1}{\delta}\int_{-\delta}^{\delta}
\bigl(1-\operatorname{Re}\varphi_{Y_n}(t)\bigr)\,dt
\longrightarrow0.
\]
Para justificar o limite da integral, use uma variável uniforme em
$[-\delta,\delta]$ e a convergência de variáveis limitadas, como na prova do
\Cref{thm:continuidade}.

O controle em um único limiar, por si só, não mostra que $Y_n\to0$ em
probabilidade. Ele mostra, porém, que $(Y_n)$ é condensada: para os índices
suficientemente grandes, a probabilidade acima é pequena; para os finitos
índices restantes, basta aumentar o limiar.

Considere uma subsequência convergente em distribuição,
$Y_{n_k}\Rightarrow Y$. Pelo \Cref{thm:continuidade}, a função característica
de $Y$ satisfaz $\varphi_Y(t)=1$ para $|t|\leqslant\delta$. Logo, para cada
um desses valores de $t$,
\[
\mathbb{E}\bigl[1-\cos(tY)\bigr]=1-\operatorname{Re}\varphi_Y(t)=0.
\]
Como a variável dentro da esperança é não negativa, $\cos(tY)=1$ quase
certamente. Escolha dois valores não nulos $t_1,t_2\in[-\delta,\delta]$ cuja
razão seja irracional. Os conjuntos
$\{2\pi k/t_1:k\in\mathbb{Z}\}$ e $\{2\pi k/t_2:k\in\mathbb{Z}\}$ têm apenas
$0$ em comum; portanto $Y=0$ quase certamente.

Pelo \Cref{lem:selecao-condensada}, toda subsequência de $(Y_n)$ possui uma
subsequência convergente em distribuição. O argumento anterior mostra que
seu limite tem a distribuição da constante $0$. Como na prova do
\Cref{thm:continuidade}, conclua que $Y_n\Rightarrow0$ e, portanto,
$Y_n\to0$ em probabilidade.
\end{exercise}

\begin{exercise}[Séries que convergem em distribuição]\label{exercise:serie-converge-dist}
Sejam $X_1,X_2,\dots$ independentes e $S_n=\sum_{m\leqslant n}X_m$. Mostre
que, se $(S_n)$ converge em distribuição, então $(S_n)$ converge em
probabilidade. Conclua, a partir do teorema de equivalência de P.~Lévy para
somas de variáveis independentes, que $(S_n)$ converge quase certamente.

\emph{Sugestão:} se $\varphi_{X_k}$ é a função característica de $X_k$, o incremento
$S_m-S_n=\sum_{n<k\leqslant m}X_k$, para $m>n$, é soma de variáveis
independentes, logo
\[
  \varphi_{S_m-S_n}(t)=\prod_{k=n+1}^{m}\varphi_{X_k}(t)=\frac{P_m(t)}{P_n(t)},
  \qquad P_j(t)=\prod_{k=1}^{j}\varphi_{X_k}(t).
\]
Se $S_n\Rightarrow S$, então $P_j(t)\to\varphi_S(t)$ para todo $t$; como
$\varphi_S$ é contínua e $\varphi_S(0)=1$, escolha $\delta>0$ com
$|\varphi_S(t)|>\frac12$ para $|t|\leqslant\delta$. Nesse intervalo os
denominadores não se anulam para $n$ grande, e o quociente $P_m(t)/P_n(t)$
tende a $1$ quando $m,n\to\infty$, para cada $t$ fixado. Pelo
\Cref{exercise:convergencia-para-zero}, segue que $S_m-S_n\to0$ em
probabilidade quando $m,n\to\infty$, isto é, $(S_n)$ é de Cauchy em
probabilidade. Escolha então $n_1<n_2<\cdots$ com
$\mathbb{P}\bigl(|S_{n_{k+1}}-S_{n_k}|>2^{-k}\bigr)\leqslant2^{-k}$; o lema de
Borel--Cantelli garante que, quase certamente, $|S_{n_{k+1}}-S_{n_k}|\leqslant
2^{-k}$ a partir de um certo índice, de modo que $(S_{n_k})$ converge quase
certamente. Como essa sequência também converge em probabilidade e a sequência
inteira é de Cauchy em probabilidade, ela converge em probabilidade para o mesmo
limite. A última conclusão, de que a convergência é quase certa e não apenas em
probabilidade, é o teorema de equivalência de Lévy: para somas de variáveis
independentes, a convergência em probabilidade acarreta a convergência quase
certa. Sua prova se apoia na desigualdade maximal de Lévy e não será
desenvolvida nestas notas.
\end{exercise}

%----------------------------------------------------------------------
% BIBLIOGRAFIA
%----------------------------------------------------------------------
\bibliographystyle{plain}
\bibliography{../arquivos-aux/referencias}

\end{document}
